% ==============================================================================
% CH08_ENTERPRISE_RAG.TEX — CHƯƠNG 8: KỸ NGHỆ RAG DOANH NGHIỆP & TÌM KIẾM LAI ĐA TẦNG
% ==============================================================================
% Định dạng: Sách chuyên khảo kinh điển (Classic Engineering Monograph)
% Chuẩn mực: Không dùng hộp đóng khung 4 phía, toán học 5 bước, bài tập đa tầng
% ==============================================================================

\chapter{Kỹ Nghệ RAG Doanh Nghiệp \& Tìm Kiếm Lai Đa Tầng}
\label{ch:enterprise-rag}

Trong các bản demo công nghệ hoặc ứng dụng mẫu dành cho người mới bắt đầu, xây dựng một hệ thống Tạo sinh Tăng cường Truy xuất (Retrieval-Augmented Generation -- RAG) dường như là bài toán vô cùng đơn giản: Đọc một file PDF, chia nhỏ thành các đoạn 512 ký tự, nạp vào một cơ sở dữ liệu véc-tơ mã nguồn mở, tìm kiếm 3 đoạn có độ tương đồng Cosine cao nhất, và ghép vào câu lệnh (Prompt) gửi tới LLM.

Tuy nhiên, khi mang kiến trúc ngây thơ này (\textbf{Naive RAG}) vào triển khai tại các ngân hàng, tập đoàn viễn thông, bệnh viện hay doanh nghiệp quy mô lớn, hệ thống sẽ sụp đổ gần như ngay lập tức. Người dùng phàn nàn vì mô hình trả lời sai mã số hợp đồng, bỏ sót các quy định trong bảng biểu, nhầm lẫn giữa các phiên bản tài liệu, và đưa ra các câu trả lời chứa đầy thông tin rác.

Một Kỹ sư AI cấp cao phải nhìn nhận RAG không phải là một chuỗi nối ghép API đơn giản, mà là một \textbf{Hệ thống Truy xuất Thông tin Chuyên sâu (Advanced Information Retrieval System)} kết hợp nhuần nhuyễn giữa:
\begin{enumerate}
  \item Tìm kiếm từ khóa chính xác kinh điển (\textbf{BM25}) và tìm kiếm ngữ nghĩa dày (\textbf{Dense Vector Retrieval}).
  \item Thuật toán chỉ mục đồ thị xấp xỉ (\textbf{HNSW}) và tái xếp hạng học sâu (\textbf{Cross-Encoder Reranking}).
  \item Đồ thị tri thức doanh nghiệp (\textbf{GraphRAG}) giải quyết bài toán suy luận toàn cục (\textit{Global Sensemaking}).
  \item Vòng lặp tự phản hồi sửa lỗi (\textbf{Self-RAG \& Corrective RAG}).
  \item Hệ thống đo lường chất lượng định lượng cấp độ công nghiệp (\textbf{RAG Triad \& Ragas}).
\end{enumerate}

---

\section{Sự Sụp Đổ Của Naive RAG \& Thực Tế Triển Khai Doanh Nghiệp}
\label{sec:naive-rag-failure}

Kiến trúc Naive RAG truyền thống dựa trên một giả định hết sức mong manh: \textit{``Nếu hai đoạn văn bản có ý nghĩa tương tự nhau, véc-tơ nhúng (Embedding) của chúng sẽ nằm gần nhau trong không gian đa chiều, và việc lấy Top-K đoạn gần nhất sẽ cung cấp đủ ngữ cảnh cho LLM.''}. Trong môi trường doanh nghiệp thực tế, giả định này sụp đổ vì 4 điểm nghẽn vật lý:

\paragraph{1. Hiện tượng ``Lạc lối ở giữa'' (Lost in the Middle)}:
Công trình nghiên cứu mang tính bước ngoặt của Liu et al. (Stanford University, 2023) đã chứng minh rằng: Khả năng truy hồi và sử dụng thông tin của các mô hình ngôn ngữ lớn tuân theo một đồ thị hình chữ U (U-shaped curve). LLM có xu hướng ghi nhớ cực tốt các thông tin nằm ở ngay đầu câu lệnh (Primacy effect) và ngay cuối câu lệnh (Recency effect), nhưng \textbf{hoàn toàn bỏ quên hoặc diễn giải sai các thông tin nằm ở quãng giữa của cửa sổ ngữ cảnh}. Khi Naive RAG nhồi nhét $10$ đến $20$ đoạn tài liệu vào prompt, nếu dữ kiện quan trọng nhất vô tình rơi vào vị trí thứ 7 hoặc thứ 8, mô hình sẽ bị ảo giác (Hallucination) dù tài liệu đúng đã được trích xuất thành công!

\paragraph{2. Ô nhiễm Ngữ cảnh (Context Pollution)}:
Khi tìm kiếm thuần túy theo véc-tơ với $K=5$, thường chỉ có 1 hoặc 2 đoạn văn bản thực sự chứa câu trả lời. 3 đoạn còn lại là những đoạn văn bản có vẻ liên quan về mặt từ ngữ nhưng chứa dữ kiện sai lệch hoặc lạc đề. Việc đưa những đoạn văn bản rác này vào prompt làm loãng tín hiệu chú ý (Attention dilution), khiến mô hình tự tin đưa ra câu trả lời sai.

\paragraph{3. Bất đối xứng Ngữ nghĩa (Query-Document Asymmetry)}:
Người dùng thường chỉ nhập một câu hỏi ngắn ngủi từ 5 đến 10 từ (ví dụ: \code{"điều kiện bồi thường bảo hiểm cháy nổ?"}), trong khi tài liệu hợp đồng là một điều khoản pháp lý phức tạp dài 500 từ. Véc-tơ biểu diễn của một câu hỏi ngắn và một điều khoản dài nằm ở hai vùng không gian rất xa nhau, khiến phép đo khoảng cách Cosine bị lệch hướng nghiêm trọng.

\paragraph{4. Sự Bất Lực Tuyệt Đối Trước Từ Khóa Chính Xác (Exact Keyword Failure)}:
Mô hình nhúng véc-tơ (Embedding Models) được huấn luyện để nắm bắt ý niệm ngữ nghĩa chung (Semantic Gist), nhưng chúng hoàn toàn mù tịt trước các ký hiệu số học và mã định danh chính xác:
\begin{itemize}
  \item Mã định danh sản phẩm: \code{SKU-X100-PRO-V2} vs \code{SKU-X100-STD-V1}.
  \item Số hợp đồng kinh tế: \code{HD-2024/098/TC} vs \code{HD-2023/098/TC}.
  \item Mã lỗi máy chủ: \code{ERR_CONNECTION_TIMED_OUT_504}.
\end{itemize}
Đối với mô hình Embedding, các chuỗi ký tự trên có khoảng cách Cosine gần như bằng nhau. Nếu người dùng hỏi về quy chế áp dụng cho hợp đồng \code{HD-2024/098/TC}, tìm kiếm véc-tơ có thể thản nhiên trả về hợp đồng năm 2023, dẫn tới hậu quả pháp lý nghiêm trọng!

---

\section{Giải Phẫu Thuật Toán Tìm Kiếm Véc-Tơ: IVFFlat vs HNSW}
\label{sec:vector-search-hnsw-ivfflat}

Để hiểu cách các cơ sở dữ liệu véc-tơ (Vector Databases như Milvus, Qdrant, Pinecone, pgvector) hoạt động, Kỹ sư AI cần nắm vững cơ sở toán học của hai cấu trúc chỉ mục tìm kiếm láng giềng gần nhất xấp xỉ (Approximate Nearest Neighbors -- ANN):

\subsection{Không Gian Metric \& Độ Đo Khoảng Cách}
Cho hai véc-tơ biểu diễn đa chiều $\mathbf{u}, \mathbf{v} \in \mathbb{R}^d$:
\begin{itemize}
  \item \textbf{Khoảng cách Euclid ($L_2$ Distance)}:
  \begin{equation}
  D_{L2}(\mathbf{u}, \mathbf{v}) = \|\mathbf{u} - \mathbf{v}\|_2 = \sqrt{\sum_{i=1}^d (u_i - v_i)^2}
  \end{equation}
  \item \textbf{Tích vô hướng (Inner Product / Dot Product)}:
  \begin{equation}
  \langle \mathbf{u}, \mathbf{v} \rangle = \mathbf{u} \cdot \mathbf{v} = \sum_{i=1}^d u_i v_i
  \end{equation}
  \item \textbf{Khoảng cách Cosine (Cosine Distance)}:
  \begin{equation}
  D_{\text{Cosine}}(\mathbf{u}, \mathbf{v}) = 1 - \frac{\mathbf{u} \cdot \mathbf{v}}{\|\mathbf{u}\|_2 \|\mathbf{v}\|_2}
  \end{equation}
\end{itemize}

\paragraph{Tại Sao Phải Chuẩn Hóa L2 Trước Khi Lưu Vào Vector DB?}:
Nếu tất cả các véc-tơ trong cơ sở dữ liệu đều được chuẩn hóa độ dài về $1$ ($\|\mathbf{u}\|_2 = \|\mathbf{v}\|_2 = 1$ ngay khi nạp vào), ta có mối liên hệ toán học hoàn hảo:
\begin{equation}
D_{L2}^2(\mathbf{u}, \mathbf{v}) = \|\mathbf{u} - \mathbf{v}\|_2^2 = \|\mathbf{u}\|_2^2 + \|\mathbf{v}\|_2^2 - 2 \mathbf{u} \cdot \mathbf{v} = 2 - 2 \langle \mathbf{u}, \mathbf{v} \rangle = 2 D_{\text{Cosine}}(\mathbf{u}, \mathbf{v})
\end{equation}
Do đó, việc tối đa hóa Cosine Similarity hoàn toàn tương đương với việc tối đa hóa Dot Product và tối thiểu hóa khoảng cách Euclid. Trong sản xuất, các kỹ sư luôn chuẩn hóa $L_2$ trước để chỉ cần tính Dot Product, tận dụng tối đa tập lệnh SIMD/AVX-512 trên CPU hoặc Tensor Cores trên GPU.

\subsection{Thuật Toán IVFFlat (Inverted File with Flat Vectors)}
IVFFlat giải quyết bài toán tìm kiếm ANN bằng phương pháp phân vùng không gian Voronoi:
\begin{enumerate}
  \item \textbf{Pha huấn luyện (Training Phase)}: Áp dụng thuật toán phân cụm K-Means lên tập dữ liệu véc-tơ để xác định $C$ điểm tâm cụm (Centroids): $\{\mathbf{c}_1, \mathbf{c}_2, \dots, \mathbf{c}_C\}$.
  \item \textbf{Xây dựng danh sách ngược (Inverted Lists)}: Mỗi véc-tơ trong cơ sở dữ liệu được gán vào tâm cụm gần nó nhất. Dữ liệu được tổ chức thành $C$ danh sách ngược (Buckets).
  \item \textbf{Pha truy vấn (Query Phase)}:
  \begin{itemize}
    \item Khi có véc-tơ truy vấn $\mathbf{q}$, hệ thống tính khoảng cách giữa $\mathbf{q}$ và tất cả $C$ tâm cụm.
    \item Chọn ra $\text{nprobe}$ cụm có tâm gần $\mathbf{q}$ nhất ($\text{nprobe} \ll C$).
    \item Quét tuần tự (Exhaustive Search) qua toàn bộ các véc-tơ nằm bên trong $\text{nprobe}$ cụm này để tìm ra Top-K kết quả cuối cùng.
  \end{itemize}
\end{enumerate}

\subsection{Thuật Toán HNSW (Hierarchical Navigable Small World)}
Được phát minh bởi Malkov và Yashunin (2018), \textbf{HNSW} là chuẩn mực vàng thống trị hầu hết các Vector Database hiện đại nhờ tốc độ tìm kiếm $O(\log N)$ vượt trội.

HNSW lấy cảm hứng từ cấu trúc dữ liệu \textbf{Skip-List} (Danh sách nhảy) nhưng mở rộng lên không gian đồ thị đa chiều:
\begin{itemize}
  \item \textbf{Cấu trúc đồ thị phân tầng}: HNSW xây dựng một hệ thống phân cấp gồm nhiều tầng đồ thị (Layer $0$ đến Layer $l_{\max}$).
  \begin{itemize}
    \item \textbf{Tầng cao nhất ($l_{\max}$)}: Chỉ chứa một số lượng rất ít các điểm dữ liệu. Đồ thị ở tầng này rất thưa (Sparse), các cạnh nối có độ dài liên kết rất lớn, cho phép thuật toán nhảy những bước khổng lồ xuyên qua không gian dữ liệu.
    \item \textbf{Tầng dưới cùng (Layer 0)}: Chứa toàn bộ $100\%$ các điểm dữ liệu trong cơ sở dữ liệu. Đồ thị ở tầng này dày đặc (Dense), phục vụ cho việc dò tìm cục bộ (Local Search) cực kỳ chi tiết.
  \end{itemize}
  \item \textbf{Cơ chế tìm kiếm}: Bắt đầu từ điểm vào (Entry Point) ở tầng đỉnh $l_{\max}$, thuật toán thực hiện tìm kiếm tham lam (Greedy Search) để tìm ra điểm gần truy vấn $\mathbf{q}$ nhất trên tầng đó. Điểm này sau đó được dùng làm điểm xuất phát để rơi xuống tầng tiếp theo. Quá trình lặp lại cho đến Layer 0.
\end{itemize}

---

\section{Tìm Kiếm Lai Đa Tầng (Multi-Stage Hybrid Search) \& RRF}
\label{sec:hybrid-search-rrf}

Để khắc phục hoàn toàn nhược điểm của tìm kiếm véc-tơ thuần túy, kiến trúc công nghiệp áp dụng mô hình \textbf{Tìm Kiếm Lai Đa Tầng (Multi-Stage Hybrid Search)}:
\begin{enumerate}
  \item \textbf{Sparse Retrieval (BM25)}: Bắt chính xác từng từ khóa, mã số hợp đồng, tên riêng.
  \item \textbf{Dense Retrieval (Vector ANN)}: Nắm bắt ngữ nghĩa, khái niệm tương đương và văn cảnh tổng quát.
  \item \textbf{Dung hợp thứ hạng (Rank Fusion)}: Kết hợp kết quả từ hai nguồn độc lập mà không cần chuẩn hóa điểm số dị biệt thông qua thuật toán \textbf{Reciprocal Rank Fusion (RRF)}.
  \item \textbf{Cross-Encoder Reranking}: Tái xếp hạng Top-50 tài liệu thu hồi bằng mô hình Transformer tương tác sâu để chọn ra Top-5 tinh túy nhất gửi cho LLM.
\end{enumerate}

\subsection{Toán Học Thuật Toán Okapi BM25}
BM25 (Best Matching 25) là đỉnh cao của mô hình xác suất trong thu hồi thông tin (Information Retrieval):
\begin{definition}[Công Thức Chấm Điểm Okapi BM25]
Cho câu truy vấn $Q = \{q_1, q_2, \dots, q_n\}$ và một tài liệu $D$ trong tập văn bản gồm $N$ tài liệu với độ dài trung bình $\text{avgdl}$:
\begin{equation}
\text{Score}_{\text{BM25}}(D, Q) = \sum_{i=1}^n \text{IDF}(q_i) \cdot \frac{f(q_i, D) \cdot (k_1 + 1)}{f(q_i, D) + k_1 \cdot \left(1 - b + b \cdot \frac{|D|}{\text{avgdl}}\right)}
\end{equation}
trong đó:
\begin{itemize}
  \item $f(q_i, D)$: Tần số xuất hiện của từ $q_i$ trong tài liệu $D$.
  \item $|D|$: Độ dài (tổng số từ) của tài liệu $D$.
  \item $k_1$: Tham số kiểm soát độ bão hòa tần số từ (mặc định $k_1 \in [1.2, 2.0]$, thường chọn $1.5$).
  \item $b$: Tham số phạt độ dài tài liệu (mặc định $b \in [0.75, 1.0]$, thường chọn $0.75$).
  \item $\text{IDF}(q_i)$: Trọng số nghịch đảo tần số tài liệu có làm mịn:
  \begin{equation}
  \text{IDF}(q_i) = \ln \left( \frac{N - n(q_i) + 0.5}{n(q_i) + 0.5} + 1 \right)
  \end{equation}
  với $n(q_i)$ là số lượng tài liệu trong tập văn bản có chứa từ $q_i$.
\end{itemize}
\end{definition}

\subsection{Dung Hợp Thứ Hạng Bằng Reciprocal Rank Fusion (RRF)}
Một thách thức lớn trong tìm kiếm lai là: Điểm số của BM25 (chạy từ $0$ đến $+\infty$, phụ thuộc độ dài văn bản) và điểm số của Dense Vector (Cosine Similarity chạy từ $-1$ đến $1$) không cùng một phân phối xác suất. Việc cộng trực tiếp $\alpha \cdot \text{Score}_{\text{BM25}} + \beta \cdot \text{Score}_{\text{Vector}}$ đòi hỏi cân chỉnh tham số rất nhạy cảm và dễ vỡ khi dữ liệu thay đổi.

Thuật toán \term{Reciprocal Rank Fusion (RRF)} do Cormack et al. đề xuất giải quyết triệt để vấn đề này bằng cách chỉ dựa vào \textbf{thứ hạng (Rank)} thay vì điểm số thô.

\begin{definition}[Công Thức Reciprocal Rank Fusion]
Cho $M$ danh sách xếp hạng tài liệu độc lập từ các hệ thống tìm kiếm khác nhau. Với mỗi tài liệu $d$, điểm số RRF tổng hợp được tính theo:
\begin{equation}
\text{RRF\_Score}(d) = \sum_{m=1}^M \frac{1}{k + r_m(d)}
\end{equation}
trong đó $r_m(d) \in \{1, 2, 3, \dots\}$ là thứ hạng (1-indexed) của tài liệu $d$ trong danh sách thứ $m$. Hằng số làm mịn $k$ (thường chọn $k = 60$) có tác dụng giảm độ dốc giữa các thứ hạng đầu, ngăn chặn việc một tài liệu đứng hạng 1 ở một danh sách đơn lẻ lấn át hoàn toàn một tài liệu đứng hạng 2 và hạng 3 ở cả hai danh sách. Nếu tài liệu $d$ không xuất hiện trong danh sách $m$, ta quy ước $r_m(d) = \infty \implies \frac{1}{k + r_m(d)} = 0$.
\end{definition}

\begin{example}[Tính Toán Bằng Tay Dung Hợp Thứ Hạng RRF]
Xét 4 tài liệu $\{A, B, C, D\}$ được xếp hạng bởi 2 động cơ:
\begin{itemize}
  \item Danh sách 1 (BM25): Hạng 1 là $A$, Hạng 2 là $B$, Hạng 3 là $C$. ($D$ không có trong danh sách).
  \item Danh sách 2 (Dense Vector): Hạng 1 là $B$, Hạng 2 là $D$, Hạng 3 là $A$. ($C$ không có trong danh sách).
\end{itemize}
Áp dụng RRF với $k = 60$:
\begin{itemize}
  \item Tài liệu $A$: $\text{RRF}(A) = \frac{1}{60 + 1} + \frac{1}{60 + 3} = \frac{1}{61} + \frac{1}{63} \approx 0.016393 + 0.015873 = 0.032266$
  \item Tài liệu $B$: $\text{RRF}(B) = \frac{1}{60 + 2} + \frac{1}{60 + 1} = \frac{1}{62} + \frac{1}{61} \approx 0.016129 + 0.016393 = 0.032522$
  \item Tài liệu $C$: $\text{RRF}(C) = \frac{1}{60 + 3} + 0 = \frac{1}{63} \approx 0.015873$
  \item Tài liệu $D$: $\text{RRF}(D) = 0 + \frac{1}{60 + 2} = \frac{1}{62} \approx 0.016129$
\end{itemize}
Thứ hạng chung cuộc: \textbf{$B > A > D > C$}.
Tài liệu $B$ giành chiến thắng vì nó giữ vị trí rất cao trên cả hai kênh tìm kiếm, chứng minh sức mạnh của nguyên lý xếp hạng tương hỗ!
\end{example}

---

\section{ColBERT: Tương Tác Muộn (Late Interaction) \& Toán Học Chấm Điểm MaxSim}
\label{sec:colbert-late-interaction}

Trong hệ thống tìm kiếm ngữ nghĩa, hai thái cực kiến trúc kinh điển đều phải đối mặt với sự đánh đổi nghiệt ngã:
\begin{itemize}
  \item \textbf{Bi-Encoder (Mô hình Hai Tháp -- Dense Embeddings)}: Nén toàn bộ ngữ cảnh của một tài liệu dài $512$ tokens thành một véc-tơ cố định duy nhất $\mathbf{e}_d \in \mathbb{R}^d$ ($d=768$ hoặc $1536$). Hiện tượng thắt cổ chai thông tin (\term{Information Bottleneck}) này làm mất mát hoàn toàn các mối liên kết cú pháp vi mô và các từ khóa hiếm. Dù tốc độ tìm kiếm cực nhanh thông qua HNSW ($O(\log N)$), độ chính xác thường xuyên bị suy giảm.
  \item \textbf{Cross-Encoder (Mô hình Tương Tác Sớm -- Early Interaction)}: Nối chuỗi câu hỏi và tài liệu $[q; d]$ để đưa qua toàn bộ các tầng Transformer với ma trận Self-Attention đầy đủ $O((|q| + |d|)^2)$. Mô hình nắm bắt trọn vẹn từng sắc thái tương tác giữa các từ, đạt độ chính xác tối thượng (Gold Standard). Tuy nhiên, chi phí tính toán bùng nổ khiến Cross-Encoder hoàn toàn bất khả thi cho việc tìm kiếm trực tiếp trên hàng triệu văn bản, mà chỉ có thể đóng vai trò tái xếp hạng (Reranker) cho vài chục ứng viên ở tầng cuối.
\end{itemize}

Kiến trúc \term{ColBERT} (Contextualized Late Interaction over BERT), do Khattab và Zaharia (Stanford University, 2020) đề xuất, đã tạo ra một cuộc cách mạng khi kết hợp hoàn hảo \textbf{tốc độ tìm kiếm của Bi-Encoder} và \textbf{độ chính xác sâu của Cross-Encoder}.

\subsection{Cơ Chế Biểu Diễn Đa Véc-Tơ \& Toán Tử MaxSim}
Thay vì nén tài liệu thành một véc-tơ đơn lẻ, ColBERT duy trì một ma trận các véc-tơ biểu diễn ở cấp độ từng token:
\begin{itemize}
  \item Câu truy vấn $q$ gồm $|q|$ tokens được ánh xạ thành ma trận:
  \begin{equation}
  E_q = \text{Normalize}\left( \text{Linear}\left( \text{BERT}(q) \right) \right) \in \mathbb{R}^{|q| \times d}
  \end{equation}
  \item Tài liệu $d$ gồm $|d|$ tokens được mã hóa ngoại tuyến thành ma trận:
  \begin{equation}
  E_d = \text{Normalize}\left( \text{Linear}\left( \text{BERT}(d) \right) \right) \in \mathbb{R}^{|d| \times d}
  \end{equation}
\end{itemize}
trong đó mỗi véc-tơ token $\mathbf{e} \in \mathbb{R}^d$ ($d$ thường chọn nhỏ, ví dụ $d=128$) đều được chuẩn hóa độ dài Euclid $\|\mathbf{e}\|_2 = 1$.

\begin{definition}[Toán Tử Chấm Điểm Tương Tác Muộn MaxSim]
Điểm số tương đồng ngữ nghĩa giữa câu truy vấn $q$ và tài liệu $d$ trong ColBERT là tổng các giá trị tích vô hướng cực đại của từng token truy vấn đối với toàn bộ các token của tài liệu:
\begin{equation}
\text{MaxSim}(q, d) = \sum_{i=1}^{|q|} \max_{j=1}^{|d|} \left( E_{q, i} \cdot E_{d, j}^T \right)
\end{equation}
\end{definition}

Về mặt hình học, toán tử MaxSim thực hiện một phép ``so khớp mềm'' (Soft Token Alignment): Với mỗi từ trong câu hỏi của người dùng, hệ thống quét qua toàn bộ tài liệu để tìm từ hoặc cụm từ có độ tương đồng ngữ nghĩa cao nhất, ghi nhận điểm số cực đại đó, rồi lấy tổng tích lũy. Nhờ đó:
\begin{enumerate}
  \item Nếu câu hỏi có 5 từ khóa quan trọng, tài liệu chỉ đạt điểm tối đa khi tìm được điểm neo thích hợp cho cả 5 từ khóa đó, loại bỏ hoàn toàn hiện tượng ảo giác do trùng lặp một từ đơn lẻ.
  \item Ma trận biểu diễn của tài liệu $E_d$ được tính toán trước ngoại tuyến (Offline Pre-computation) và lưu trữ trong chỉ mục, loại bỏ hoàn toàn nhu cầu chạy lại Transformer cho tài liệu tại thời điểm người dùng truy vấn.
\end{enumerate}

\subsection{Bảng So Sánh Toàn Diện Các Mô Hình Truy Xuất}
\begin{table}[htbp]
\centering
\small
\begin{tabular}{lccc}
\toprule
\textbf{Tiêu chí Đánh giá} & \textbf{Bi-Encoder (Dense)} & \textbf{Cross-Encoder} & \textbf{ColBERT (Late Interaction)} \\
\midrule
Biểu diễn tài liệu & $1$ véc-tơ duy nhất & Ghép cặp trực tiếp $[q; d]$ & Ma trận đa véc-tơ $|d| \times 128$ \\
Tiền tính toán tài liệu & Có (Offline Indexing) & Không thể & Có (Offline Indexing) \\
Độ phức tạp tính điểm & $O(d)$ & $O((|q|+|d|)^2 \cdot L)$ & $O(|q| \cdot |d| \cdot d)$ (Tích ma trận) \\
Dung lượng lưu trữ RAM & Thấp ($\approx 1.5\text{ KB/doc}$) & $0$ (Không lưu chỉ mục) & Cao ($\approx 100\text{ KB/doc}$) \\
Độ trễ truy vấn & $< 10\text{ ms}$ (HNSW) & $> 500\text{ ms}$ (Heavy GPU) & $< 30\text{ ms}$ (PLAID Index) \\
Độ chính xác truy xuất & Trung bình - Khá & Xuất sắc (Chuẩn mực) & Tiệm cận Cross-Encoder \\
\bottomrule
\end{tabular}
\caption{So sánh đặc tính kỹ thuật giữa Bi-Encoder, Cross-Encoder và ColBERT}
\label{tab:retrieval-comparison}
\end{table}

\subsection{Đột Phá Đa Phương Thức: ColPali \& Loại Bỏ Quy Trình OCR}
Một bước tiến vượt bậc của cơ chế Late Interaction là sự ra đời của \term{ColPali} (Faysse et al., 2024). Trong các ứng dụng RAG truyền thống cho tài liệu PDF, hóa đơn hay báo cáo tài chính, kỹ sư bắt buộc phải sử dụng các pipeline OCR phức tạp (như Tesseract, PaddleOCR) để trích xuất văn bản thô, dẫn tới việc làm vỡ nát các cấu trúc bảng biểu, biểu đồ và hình ảnh minh họa.

ColPali tích hợp trực tiếp mô hình Ngôn ngữ - Thị giác (Vision LLM như PaliGemma):
\begin{enumerate}
  \item Mỗi trang PDF được render thành một hình ảnh độ phân giải cao và chia thành một lưới các mảnh ảnh cục bộ (Image Patches, ví dụ $32 \times 32$ patches).
  \item Vision Transformer mã hóa mỗi patch thành một véc-tơ đa chiều, tạo thành ma trận biểu diễn $E_{\text{page}} \in \mathbb{R}^{1024 \times 128}$.
  \item Câu hỏi dạng văn bản của người dùng được mã hóa thành các véc-tơ token $E_q$.
  \item Điểm số phù hợp của trang tài liệu được tính toán trực tiếp bằng toán tử MaxSim giữa các token câu hỏi và các patch hình ảnh.
\end{enumerate}
Cơ chế này cho phép hệ thống tìm kiếm chuẩn xác từng con số nằm sâu trong góc biểu đồ hoặc bảng biểu tài chính mà không cần chạy bất kỳ một dòng lệnh OCR nào!

\subsection{Tối Ưu Hóa Tốc Độ Bằng Chỉ Mục PLAID}
Thách thức lớn nhất của ColBERT là chi phí lưu trữ: Với 1 triệu tài liệu, mỗi tài liệu 200 tokens, hệ thống phải lưu trữ tới 200 triệu véc-tơ! Để giải quyết bài toán này, Santhanam et al. (2022) giới thiệu cấu trúc chỉ mục \term{PLAID} (Performance-optimized Late Interaction with Approximated Information and Distillation):
\begin{itemize}
  \item \textbf{Phân cụm tâm véc-tơ (Centroid Clustering)}: Nhóm các véc-tơ token vào các cụm $K$-Means trong không gian đa chiều.
  \item \textbf{Tỉa nhánh ứng viên nhanh}: Khi có truy vấn, chỉ quét qua các tài liệu có chứa các token thuộc cùng tâm cụm với các token truy vấn.
  \item \textbf{Lượng tử hóa phần dư (Residual Quantization)}: Nén mỗi véc-tơ token từ 128 số thực FP32 (512 bytes) xuống chỉ còn 16 đến 32 bytes bằng cách chỉ lưu mã cụm và sai số lượng tử hóa. Nhờ đó, dung lượng chỉ mục giảm hơn $16\times$ lần, đưa thời gian truy vấn về dưới $20\text{ ms}$.
\end{itemize}

---

\section{GraphRAG: Đồ Thị Tri Thức, Phân Cụm Leiden \& Suy Luận Toàn Cầu}
\label{sec:graph-rag}

\subsection{Giới Hạn Của Vector RAG Trước Câu Hỏi Toàn Cục}
Một trong những điểm mù lớn nhất của Vector RAG là bài toán \term{Global Sensemaking} (Suy luận toàn cục).
Giả sử doanh nghiệp có 5.000 báo cáo kiểm toán nội bộ. Người dùng đặt câu hỏi:
\begin{center}
\textit{``Các rủi ro tuân thủ pháp lý phổ biến nhất xuyên suốt toàn bộ các chi nhánh năm qua là gì?''}
\end{center}
Vector RAG hoàn toàn bất lực trước câu hỏi này! Vì sao?
Vector Search được thiết kế cho bài toán \textit{``Tìm kiếm kim trong đáy bể''} (Needle in a haystack) -- tức tìm các đoạn văn bản cụ thể khớp với một thực thể rõ ràng. Nhưng với câu hỏi trên, không có bất kỳ một đoạn văn bản (chunk) đơn lẻ nào chứa câu trả lời tổng hợp. Câu trả lời đòi hỏi phải liên kết thông tin phân mảnh từ hàng trăm tài liệu khác nhau.

\subsection{Kiến Trúc GraphRAG: Từ Trích Xuất Bộ Ba Tới Phân Cụm Cộng Đồng}
Để giải quyết bài toán này, Microsoft Research đã công bố kiến trúc \term{GraphRAG} (Edge et al., 2024) kết hợp giữa Đồ thị Tri thức (Knowledge Graph) và Mô hình Ngôn ngữ Lớn:

\begin{figure}[htbp]
\centering
\begin{tikzpicture}[node distance=1.0cm, auto,
    box/.style={rectangle, draw=primarycolor, thick, fill=boxbg, text centered, rounded corners=2pt, minimum height=0.75cm, font=\small},
    arrow/.style={draw=primarycolor, -Stealth, thick}
]
    \node[box, text width=9.5cm] (b1) {1. Extraction: LLM trích xuất Thực thể (Entities) \& Quan hệ (Relations)};
    \node[box, below=0.35cm of b1, text width=9.5cm, fill=accentcolor!20] (b2) {2. Graph Construction: Xây dựng Đồ thị Tri thức $G = (V, E)$};
    \node[box, below=0.35cm of b2, text width=9.5cm, fill=yellow!15] (b3) {3. Community Detection: Thuật toán Leiden phân cụm đồ thị phân cấp};
    \node[box, below=0.35cm of b3, text width=9.5cm, fill=green!15] (b4) {4. Community Summarization: LLM tóm tắt báo cáo cho từng cộng đồng tri thức};
    \node[box, below=0.35cm of b4, text width=9.5cm, fill=primarycolor!15] (b5) {5. Global Search: Map-Reduce các bản tóm tắt cộng đồng để sinh câu trả lời toàn cục};

    \path[arrow] (b1) -- (b2);
    \path[arrow] (b2) -- (b3);
    \path[arrow] (b3) -- (b4);
    \path[arrow] (b4) -- (b5);
\end{tikzpicture}
\caption{Quy trình lập chỉ mục và truy xuất tri thức phân cấp trong GraphRAG}
\label{fig:graph-rag-pipeline}
\end{figure}

\begin{enumerate}
  \item \textbf{Trích xuất Thực thể \& Quan hệ (Entity-Relation Extraction)}:
  Sử dụng LLM quét qua từng chunk tài liệu để trích xuất các bộ ba thực thể (Triplets):
  \begin{equation}
  (\text{Entity}_A, \text{Relationship}, \text{Entity}_B) \quad \text{ví dụ: } (\text{vLLM}, \text{implements}, \text{PagedAttention})
  \end{equation}
  \item \textbf{Phát hiện Cộng đồng Phân cấp Bằng Thuật Toán Leiden}:
  Thuật toán phân cụm đồ thị kinh điển Louvain thường gặp phải hiện tượng nghiêm trọng: tạo ra các cộng đồng bị ngắt kết nối cục bộ (Disconnected Communities) hoặc các nút cô lập bên trong cụm. GraphRAG sử dụng thuật toán \term{Leiden} (Traag et al., 2019) với bước tinh chỉnh phân hoạch con (Sub-partition refinement), đảm bảo mọi cộng đồng được phát hiện đều liên thông chặt chẽ. Đồ thị được phân tầng thành nhiều cấp độ:
  \begin{itemize}
    \item \textbf{Cấp độ 0 (C0 -- Macro)}: Vài cụm chủ đề bao quát toàn bộ kho dữ liệu.
    \item \textbf{Cấp độ 1 (C1 -- Meso)}: Các phân ngành và nhóm lĩnh vực nghiệp vụ.
    \item \textbf{Cấp độ 2 (C2 -- Micro)}: Các cụm chi tiết liên kết trực tiếp giữa các thực thể vi mô.
  \end{itemize}
  \item \textbf{Tóm tắt Cộng đồng (Community Summaries)}:
  Với mỗi cộng đồng ở từng cấp độ, LLM tạo một bản báo cáo điều hành cấu trúc (Executive Summary) gồm các phát hiện chính, mức độ tác động và các bằng chứng thực nghiệm.
  \item \textbf{Hai Chế Độ Truy Vấn Chuẩn}:
  \begin{itemize}
    \item \textbf{Global Search (Map-Reduce)}: Phục vụ câu hỏi khái quát. Hệ thống chia câu hỏi thành các tác vụ song song gửi tới các Community Summaries (pha Map) để thu thập các điểm đánh giá và bằng chứng, sau đó tổng hợp thành câu trả lời hoàn chỉnh (pha Reduce).
    \item \textbf{Local Search ($k$-hop Expansion)}: Phục vụ câu hỏi cụ thể về một thực thể. Hệ thống dùng Vector Search tìm các thực thể hạt nhân, sau đó mở rộng mạng lưới láng giềng $k$-hop để thu thập các mối quan hệ trực tiếp và văn bản gốc.
  \end{itemize}
\end{enumerate}

---

\section{Contextual Retrieval \& Tối Ưu Chi Phí Với Prompt Caching}
\label{sec:contextual-retrieval}

\subsection{Khuyết Tật Cô Lập Đoạn Văn Bản (The Chunk Isolation Problem)}
Trong các hệ thống RAG truyền thống, các tài liệu lớn luôn được chia nhỏ thành các đoạn độc lập có kích thước cố định (ví dụ $300 - 500$ tokens). Việc bẻ gãy văn bản này tạo ra một vấn đề chí tử: \term{Chunk Isolation} (Sự cô lập ngữ cảnh).

Hãy xem xét đoạn văn bản sau trích từ một báo cáo tài chính 100 trang:
\begin{quote}\itshape
``Doanh thu thuần trong quý này đạt 45.000 tỷ VND, tăng trưởng 25\% so với cùng kỳ năm ngoái nhờ sự bùng nổ của mảng xe điện. Tuy nhiên, biên lợi nhuận gộp giảm nhẹ 1.2\% do chi phí đầu vào tăng cao.''
\end{quote}
Đoạn văn bản trên hoàn toàn không nhắc tới: Đây là công ty nào? Báo cáo thuộc năm tài chính nào? Quý mấy?
Khi người dùng đặt câu hỏi: \textit{``Doanh thu quý 3 năm 2024 của VinFast là bao nhiêu?''}, cả thuật toán tìm kiếm từ khóa BM25 lẫn tìm kiếm véc-tơ Dense Embeddings đều \textbf{hoàn toàn bỏ qua} đoạn văn bản trên vì trong đoạn không hề xuất hiện từ khóa \code{"VinFast"} hay \code{"2024"}!

\subsection{Cơ Chế Bổ Sung Ngữ Cảnh Tự Động Của Anthropic}
Để giải quyết triệt để vấn đề này, phòng nghiên cứu Anthropic (2024) đã công bố kỹ thuật \term{Contextual Retrieval}. Trước khi nạp các chunk vào cơ sở dữ liệu, một mô hình ngôn ngữ nhỏ được yêu cầu đọc toàn bộ tài liệu gốc và sinh ra một đoạn giải thích ngắn gọn (từ $50$ đến $100$ tokens) đặt ngữ cảnh cho riêng chunk đó:

\begin{figure}[htbp]
\centering
\begin{tikzpicture}[node distance=1.0cm, auto,
    raw_box/.style={rectangle, draw=secondarycolor, dashed, fill=accentcolor!10, text width=5.5cm, align=left, font=\scriptsize, inner sep=5pt},
    ctx_box/.style={rectangle, draw=primarycolor, thick, fill=boxbg, text width=5.5cm, align=left, font=\scriptsize, inner sep=5pt},
    arrow/.style={draw=primarycolor, -Stealth, thick}
]
    \node[raw_box] (raw) {\textbf{[Chunk Thô Gốc]}:\\``Doanh thu thuần trong quý này đạt 45.000 tỷ VND, tăng trưởng 25\%...''};
    \node[ctx_box, right=1.2cm of raw] (aug) {\textbf{[Contextual Chunk Hoàn Chỉnh]}:\\
    \textcolor{primarycolor}{\textbf{<document\_context>}}\\
    \textit{Đoạn trích này nằm trong Báo cáo tài chính Quý 3 năm 2024 của Tập đoàn Vingroup, phần phân tích kết quả kinh doanh của công ty thành viên VinFast.}\\
    \textcolor{primarycolor}{\textbf{</document\_context>}}\\
    \textcolor{secondarycolor}{\textbf{<chunk\_content>}}\\
    Doanh thu thuần trong quý này đạt 45.000 tỷ VND...\\
    \textcolor{secondarycolor}{\textbf{</chunk\_content>}}};

    \path[arrow] (raw) -- node[above, font=\tiny\bfseries] {Context Injection} (aug);
\end{tikzpicture}
\caption{Cấu trúc một Contextual Chunk được bổ sung ngữ cảnh tài liệu}
\label{fig:contextual-retrieval}
\end{figure}

Đoạn văn bản hoàn chỉnh sau khi ghép ngữ cảnh được nạp song song vào cả hai hệ thống:
\begin{itemize}
  \item \textbf{Chỉ mục BM25}: Bắt chính xác các từ khóa thực thể vừa được bổ sung (\code{Vingroup}, \code{VinFast}, \code{Quý 3/2024}).
  \item \textbf{Chỉ mục Dense Vector}: Véc-tơ nhúng được kéo về đúng vị trí ngữ nghĩa của doanh nghiệp và niên độ tài chính.
\end{itemize}
Thực nghiệm của Anthropic trên nhiều tập dữ liệu doanh nghiệp chứng minh: Việc áp dụng Contextual Retrieval kết hợp tìm kiếm lai BM25 + Dense giúp \textbf{giảm tới 49\% tỷ lệ truy xuất thất bại}, và khi kết hợp thêm bộ tái xếp hạng Reranker, tỷ lệ lỗi giảm kỷ lục tới \textbf{67\%}!

\subsection{Giải Quyết Rào Cản Chi Phí Bằng Prompt Caching}
Thoạt nhìn, chi phí để áp dụng Contextual Retrieval cho một kho tài liệu doanh nghiệp hàng triệu tokens dường như là bất khả thi: Mỗi chunk cần đưa toàn bộ tài liệu $100$ trang ($50.000$ tokens) vào LLM để sinh $50$ tokens ngữ cảnh! Nếu một tài liệu có $100$ chunks, ta phải gửi $100 \times 50.000 = 5.000.000$ tokens đầu vào!

Tuy nhiên, đây chính là lúc đòn bẩy kỹ thuật của \textbf{Prompt Caching} (đã phân tích sâu tại Chương~\ref{ch:llm-runtime-tokenomics}) phát huy sức mạnh tối thượng:
Toàn bộ văn bản tài liệu 100 trang được đưa vào hệ thống như một tiền tố tĩnh cố định (\textit{Static Prefix}). Nhờ cơ chế bộ nhớ đệm KV Cache của các nhà cung cấp mô hình (Anthropic Prompt Caching, OpenAI Prefix Cache):
\begin{equation}
\text{Cost}_{\text{prefill}} = \text{Cost}_{\text{write\_cache}} + 99 \times \text{Cost}_{\text{read\_cache}} \approx 0.1 \times \text{Cost}_{\text{uncached}}
\end{equation}
Chi phí đọc lại KV Cache giảm tới $90\%$, biến giải pháp Contextual Retrieval từ một ý tưởng lý thuyết xa xỉ thành một chuẩn mực kỹ nghệ bắt buộc trong mọi dự án RAG doanh nghiệp hiện đại.

---

\section{Adaptive RAG: Corrective RAG (CRAG) \& Self-RAG}
\label{sec:adaptive-rag}

Một hệ thống RAG tĩnh chỉ thực hiện một chiều: \code{Retrieve -> Generate}. Nếu tài liệu trích xuất sai, LLM chắc chắn sinh ảo giác.
Kiến trúc thế hệ mới chuyển đổi RAG thành một quy trình có khả năng \textbf{tự nhận thức và tự sửa lỗi (Self-Correction)}:

\begin{figure}[htbp]
\centering
\begin{tikzpicture}[node distance=1.3cm, auto,
    decision/.style={diamond, draw=warncolor, thick, fill=yellow!15, align=center, font=\footnotesize, aspect=1.6},
    block/.style={rectangle, draw=primarycolor, thick, fill=boxbg, align=center, rounded corners=2pt, minimum height=0.8cm, font=\small},
    arrow/.style={draw=primarycolor, -Stealth, thick}
]
    \node[block] (query) {User Query};
    \node[block, right=0.8cm of query] (ret) {Retrieval};
    \node[decision, right=0.8cm of ret] (eval) {Retrieval\\Quality?};
    \node[block, above right=0.6cm and 0.6cm of eval, fill=green!15] (correct) {Generate Answer\\(Citations)};
    \node[block, below right=0.6cm and 0.6cm of eval, fill=red!15] (web) {Fallback Web Search\\(Tavily / Serper)};
    \node[block, right=1.0cm of eval, fill=accentcolor!25] (refine) {Knowledge Refinement\\(Strip-level filtering)};

    \path[arrow] (query) -- (ret);
    \path[arrow] (ret) -- (eval);
    \path[arrow] (eval) -- node[above, font=\tiny\bfseries] {Correct} (correct);
    \path[arrow] (eval) -- node[below, font=\tiny\bfseries] {Incorrect} (web);
    \path[arrow] (eval) -- node[above, font=\tiny\bfseries] {Ambiguous} (refine);
    \path[arrow] (refine) -- (correct);
    \path[arrow] (web) |- (correct);
\end{tikzpicture}
\caption{Luồng công việc tự thích ứng trong Corrective RAG (CRAG)}
\label{fig:crag-flowchart}
\end{figure}

\begin{itemize}
  \item \textbf{Đánh giá chất lượng trích xuất (Retrieval Evaluator)}: Một mô hình phân loại nhẹ tính toán mức độ tin cậy $\mathcal{C}$ của các tài liệu thu hồi đối với câu hỏi:
  \begin{itemize}
    \item Nếu $\mathcal{C} \ge \gamma_{\text{high}}$ (\code{Correct}): Dữ liệu hoàn toàn đáng tin cậy $\implies$ Đẩy thẳng vào LLM để sinh câu trả lời có trích dẫn.
    \item Nếu $\mathcal{C} < \gamma_{\text{low}}$ (\code{Incorrect}): Tài liệu thu hồi hoàn toàn lạc đề $\implies$ Kích hoạt tìm kiếm Internet (Web Search Fallback) hoặc từ chối trả lời an toàn.
    \item Nếu $\gamma_{\text{low}} \le \mathcal{C} < \gamma_{\text{high}}$ (\code{Ambiguous}): Tài liệu có chứa thông tin nhưng bị lẫn tạp âm $\implies$ Kích hoạt bước chia nhỏ tri thức (Knowledge Stripping) để lọc bỏ câu từ dư thừa trước khi đưa vào LLM.
  \end{itemize}
\end{itemize}

---

\section{Đo Lường Chất Lượng RAG: Khung Đánh Giá Ragas \& ARES}
\label{sec:ragas-evaluation}

Trong môi trường doanh nghiệp, bạn không thể đánh giá RAG bằng cảm tính hay một vài câu hỏi thử nghiệm ngẫu nhiên. Mọi nâng cấp (thay đổi chunk size, đổi embedding model, tinh chỉnh top-k) phải được chứng minh bằng các chỉ số định lượng.

Khung đo lường \textbf{RAG Triad} chuẩn mực bao gồm 3 chỉ số cốt lõi:

\begin{figure}[htbp]
\centering
\begin{tikzpicture}[auto,
    node_t/.style={circle, draw=primarycolor, thick, fill=boxbg, align=center, minimum size=2.0cm, font=\small\bfseries},
    edge_label/.style={font=\footnotesize\bfseries, align=center}
]
    \node[node_t] (query) at (0, 0) {Query};
    \node[node_t] (ctx) at (5.0, 0) {Context};
    \node[node_t] (ans) at (2.5, -2.8) {Answer};

    \draw[<->, draw=primarycolor, thick] (query) -- node[above, edge_label] {Context Relevance} (ctx);
    \draw[<->, draw=secondarycolor, thick] (ctx) -- node[right, edge_label, yshift=-0.2cm] {Faithfulness\\(Groundedness)} (ans);
    \draw[<->, draw=accentcolor, thick] (ans) -- node[left, edge_label, yshift=-0.2cm] {Answer Relevance} (query);
\end{tikzpicture}
\caption{Tam giác chất lượng đánh giá RAG (RAG Triad)}
\label{fig:rag-triad}
\end{figure}

\subsection{1. Độ Trung Thực (Faithfulness / Groundedness)}
Đo lường mức độ mà câu trả lời của LLM được chứng minh bởi ngữ cảnh trích xuất (chống ảo giác).
\begin{equation}
\text{Faithfulness} = \frac{|\text{Số lượng khẳng định (Claims) được chứng minh bởi Context}|}{|\text{Tổng số khẳng định (Claims) có trong Answer}|}
\end{equation}

\subsection{2. Độ Phù Hợp Của Câu Trả Lời (Answer Relevance)}
Đo lường mức độ câu trả lời giải quyết đúng trọng tâm câu hỏi của người dùng (tránh trả lời vòng vo, lạc đề).
Để đo lường khách quan mà không cần ground truth, ta dùng LLM sinh ngược $N$ câu hỏi tiềm năng từ câu trả lời đã sinh: $\{Q_1', Q_2', \dots, Q_N'\}$. Sau đó đo độ tương đồng Cosine trung bình giữa câu hỏi gốc $Q$ và các câu hỏi được sinh ngược:
\begin{equation}
\text{Answer Relevance} = \frac{1}{N} \sum_{i=1}^N \cos(\mathbf{e}_Q, \mathbf{e}_{Q_i'})
\end{equation}

\subsection{3. Độ Chính Xác Ngữ Cảnh (Context Precision)}
Đo lường xem các tài liệu thực sự liên quan có được sắp xếp ở những vị trí đầu tiên của danh sách trích xuất hay không:
\begin{equation}
\text{Context Precision@K} = \frac{\sum_{k=1}^K \text{Precision@k} \times v_k}{\text{Tổng số tài liệu liên quan trong Top-K}}
\end{equation}
trong đó $v_k \in \{0, 1\}$ là chỉ số nhị phân báo hiệu tài liệu ở hạng $k$ có liên quan hay không.

---

\section{Hiện Thực Hóa: Bộ Động Cơ Tìm Kiếm Lai Đa Tầng \& Đánh Giá}
\label{sec:python-hybrid-engine}

Dưới đây là mã nguồn Python chuẩn mực, tự chứa và không phụ thuộc vào framework bên ngoài, hiện thực hóa hoàn chỉnh thuật toán Okapi BM25, dung hợp thứ hạng RRF, tái xếp hạng Cross-Encoder hai tầng, mô-đun mở rộng đồ thị GraphRAG, và bộ đo lường Faithfulness.

\subsection{1. Okapi BM25 \& Reciprocal Rank Fusion}
\begin{lstlisting}[language=Python]
import math
from collections import Counter
from typing import List, Tuple, Dict

class SimpleBM25:
    """Hien thuc hoa thuat toan Okapi BM25 chuan xac."""
    def __init__(self, corpus: List[List[str]], k1: float = 1.5, b: float = 0.75):
        self.k1 = k1
        self.b = b
        self.corpus_size = len(corpus)
        self.doc_lens = [len(doc) for doc in corpus]
        self.avgdl = sum(self.doc_lens) / self.corpus_size if self.corpus_size > 0 else 0
        self.doc_freqs: Dict[str, int] = Counter()
        self.doc_term_counts = []

        for doc in corpus:
            term_counts = Counter(doc)
            self.doc_term_counts.append(term_counts)
            for term in term_counts:
                self.doc_freqs[term] += 1

    def idf(self, term: str) -> float:
        n_q = self.doc_freqs.get(term, 0)
        return math.log((self.corpus_size - n_q + 0.5) / (n_q + 0.5) + 1.0)

    def score(self, query: List[str], doc_idx: int) -> float:
        score = 0.0
        doc_len = self.doc_lens[doc_idx]
        term_counts = self.doc_term_counts[doc_idx]

        for term in query:
            if term not in term_counts:
                continue
            tf = term_counts[term]
            idf_val = self.idf(term)
            num = tf * (self.k1 + 1.0)
            den = tf + self.k1 * (1.0 - self.b + self.b * (doc_len / self.avgdl))
            score += idf_val * (num / den)

        return score

    def search(self, query: List[str], top_k: int = 5) -> List[Tuple[int, float]]:
        scores = [(idx, self.score(query, idx)) for idx in range(self.corpus_size)]
        scores.sort(key=lambda x: x[1], reverse=True)
        return scores[:top_k]

def reciprocal_rank_fusion(
    ranked_lists: List[List[int]],
    k: int = 60
) -> List[Tuple[int, float]]:
    """Dung hop thu hang Reciprocal Rank Fusion (RRF)."""
    scores: Dict[int, float] = {}
    for rank_list in ranked_lists:
        for rank_idx, doc_id in enumerate(rank_list):
            rank_1_indexed = rank_idx + 1
            scores[doc_id] = scores.get(doc_id, 0.0) + (1.0 / (k + rank_1_indexed))

    return sorted(scores.items(), key=lambda x: x[1], reverse=True)
\end{lstlisting}

\subsection{2. Tái Xếp Hạng Hai Tầng (Two-Stage Hybrid Retriever)}
\begin{lstlisting}[language=Python]
class TwoStageHybridRetriever:
    """Ket hop BM25 va Dense Search qua RRF, sau do Rerank bang Cross-Encoder."""
    def __init__(self, corpus: List[List[str]], bm25_engine: SimpleBM25, cross_encoder_fn):
        self.corpus = corpus
        self.bm25 = bm25_engine
        self.reranker = cross_encoder_fn

    def retrieve(self, query: List[str], dense_ranked_ids: List[int], 
                 top_k_stage1: int = 10, top_k_final: int = 3) -> List[Tuple[int, float]]:
        # 1. Sparse Search qua BM25
        bm25_results = self.bm25.search(query, top_k=top_k_stage1)
        sparse_ranked = [doc_id for doc_id, _ in bm25_results]

        # 2. Stage 1 Fusion qua RRF
        rrf_ranked = reciprocal_rank_fusion([sparse_ranked, dense_ranked_ids], k=60)
        candidate_ids = [doc_id for doc_id, _ in rrf_ranked[:top_k_stage1]]

        # 3. Stage 2 Rerank chuyen sau bang Cross-Encoder
        query_str = " ".join(query)
        scored_candidates = []
        for doc_id in candidate_ids:
            doc_text = " ".join(self.corpus[doc_id])
            score = self.reranker(query_str, doc_text)
            scored_candidates.append((doc_id, score))

        scored_candidates.sort(key=lambda x: x[1], reverse=True)
        return scored_candidates[:top_k_final]
\end{lstlisting}

\subsection{3. Module Đánh Giá Faithfulness \& Mở Rộng Đồ Thị GraphRAG}
\begin{lstlisting}[language=Python]
class FaithfulnessEvaluator:
    """Danh gia ty le khang dinh trong cau tra loi duoc chung minh boi ngu canh."""
    def evaluate(self, claims: List[str], context: str) -> Dict[str, float]:
        if not claims:
            return {"faithfulness": 1.0, "supported_count": 0, "total_claims": 0}

        context_lower = context.lower()
        supported = 0
        for claim in claims:
            claim_tokens = [w for w in claim.lower().split() if len(w) > 2]
            if all(token in context_lower for token in claim_tokens):
                supported += 1

        score = supported / len(claims)
        return {
            "faithfulness": score,
            "supported_count": supported,
            "total_claims": len(claims)
        }

class GraphEntityRetriever:
    """Truy xuat mo rong lang gieng k-hop tren do thi tri thuc."""
    def __init__(self, adjacency_list: Dict[str, List[str]]):
        self.adj = adjacency_list

    def expand_subgraph(self, seed_entities: List[str], k_hops: int = 1) -> set:
        visited = set(seed_entities)
        current_frontier = set(seed_entities)

        for _ in range(k_hops):
            next_frontier = set()
            for entity in current_frontier:
                neighbors = self.adj.get(entity, [])
                for neighbor in neighbors:
                    if neighbor not in visited:
                        visited.add(neighbor)
                        next_frontier.add(neighbor)
            current_frontier = next_frontier

        return visited
\end{lstlisting}

\subsection{4. Động Cơ Chấm Điểm ColBERT MaxSim}
\begin{lstlisting}[language=Python]
import numpy as np

class ColBERTMaxSimEngine:
    """
    Mo phong cach cham diem Late Interaction trong ColBERT (Khattab & Zaharia, 2020).
    Score(Q, D) = sum_{i=1}^{|Q|} max_{j=1}^{|D|} (E_q[i] @ E_d[j].T)
    Bao toan su can chinh ngu nghia o cap do token thay vi nen ve mot vector duy nhat.
    """
    @staticmethod
    def compute_maxsim(query_embeddings: np.ndarray, doc_embeddings: np.ndarray) -> float:
        # Similarity matrix S: (n_q_tokens, n_d_tokens)
        sim_matrix = np.matmul(query_embeddings, doc_embeddings.T)
        # Max per query token along document tokens dimension
        max_per_query_token = np.max(sim_matrix, axis=1)
        # Sum across all query tokens
        return float(np.sum(max_per_query_token))

    @staticmethod
    def rank_documents(
        query_embeddings: np.ndarray,
        corpus_embeddings: List[np.ndarray]
    ) -> List[Tuple[int, float]]:
        scores = [
            (idx, ColBERTMaxSimEngine.compute_maxsim(query_embeddings, doc_emb))
            for idx, doc_emb in enumerate(corpus_embeddings)
        ]
        scores.sort(key=lambda x: x[1], reverse=True)
        return scores
\end{lstlisting}

\subsection{5. Bộ Tiền Xử Lý Ngữ Cảnh Hóa Contextual Retrieval}
\begin{lstlisting}[language=Python]
class ContextualChunkAugmenter:
    """
    Mo phong Contextual Retrieval cua Anthropic (2024).
    Bo sung ngu canh tong the cua tai lieu vao tung chunk truoc khi lap chi muc.
    """
    def __init__(self, doc_context_map: Dict[str, str]):
        self.doc_context_map = doc_context_map

    def augment_chunk(self, doc_id: str, raw_chunk: str) -> str:
        ctx = self.doc_context_map.get(doc_id, "")
        if not ctx:
            return raw_chunk
        return f"<document_context>\n{ctx}\n</document_context>\n\n<chunk_content>\n{raw_chunk}\n</chunk_content>"

    def build_hybrid_search_corpus(
        self, raw_chunks: List[Tuple[str, str]]
    ) -> List[str]:
        return [self.augment_chunk(doc_id, chunk) for doc_id, chunk in raw_chunks]
\end{lstlisting}

\begin{simulationbox}[Mô Phỏng Trực Quan Hóa ColBERT MaxSim \& GraphRAG]
Dưới đây là các kịch bản chuẩn để yêu cầu AI sinh mã nguồn hoạt họa 60 FPS bằng thư viện Manim Community Edition (Manim CE):

\paragraph{Prompt 1: Hoạt Họa 60 FPS Cơ Chế So Khớp Tương Tác Muộn ColBERT MaxSim}
\begin{quote}\small\itshape
``Hãy viết một đoạn mã Python hoàn chỉnh sử dụng thư viện Manim CE (\code{from manim import *}) để tạo một video giải thích trực quan về cơ chế ColBERT Late Interaction và ma trận MaxSim với các yêu cầu:
\begin{enumerate}
  \item \textbf{Cảnh 1: Tách Biệt Biểu Diễn Token}: Hiển thị câu hỏi $Q$ gồm 3 tokens màu xanh lam nhạt và tài liệu $D$ gồm 6 tokens màu vàng cam. Từng token được phóng chiếu thành các véc-tơ đa chiều riêng biệt thay vì bị nén phẳng thành 1 véc-tơ đơn.
  \item \textbf{Cảnh 2: Ma Trận Tương Đồng $3 \times 6$}: Thể hiện ma trận tích vô hướng $S = E_q \cdot E_d^T$ kích thước $3 \times 6$ sáng lên, với các ô có giá trị tương đồng cao mang sắc thái màu xanh lục rực rỡ.
  \item \textbf{Cảnh 3: Quét Tìm Max Trên Từng Dòng}: Một con trỏ đỏ quét qua từng dòng (tương ứng với từng token câu hỏi), tự động đóng khung giá trị lớn nhất ($\max$) trên dòng đó.
  \item \textbf{Cảnh 4: Cộng Dồn Tổng Điểm MaxSim}: Ba giá trị cực đại được trích xuất và kéo xuống một thanh tính tổng ($\sum$), cộng lại thành điểm số tương đồng MaxSim chung cuộc.
  \item Render trực tiếp bằng lệnh: \code{manim -pql scene.py ColBERTMaxSimScene}.''
\end{enumerate}
\end{quote}

\paragraph{Prompt 2: Hoạt Họa 60 FPS GraphRAG: Phân Cụm Leiden \& Map-Reduce Global Search}
\begin{quote}\small\itshape
``Hãy viết một kịch bản Manim CE 60 FPS mô phỏng kiến trúc GraphRAG từ phân cụm đồ thị tri thức đến truy vấn toàn cục:
\begin{enumerate}
  \item \textbf{Cảnh 1: Mạng Lưới Đồ Thị Tri Thức}: Xuất hiện 30 nút thực thể (Entities) và các cạnh quan hệ (Relations) chằng chit trong không gian 2D.
  \item \textbf{Cảnh 2: Phân Cụm Phân Cấp Leiden}: Thuật toán Leiden kích hoạt một làn sóng quét qua đồ thị, tự động gom các thực thể thành 4 cộng đồng màu sắc độc lập (C0 vĩ mô). Camera phóng to vào từng cụm thể hiện việc phân tách tiếp thành các tiểu cụm (C1 và C2).
  \item \textbf{Cảnh 3: Sinh Báo Cáo Tóm Tắt (Community Summaries)}: Mỗi cụm cô đọng lại thành một cuốn sổ báo cáo điều hành vĩ mô.
  \item \textbf{Cảnh 4: Truy Vấn Toàn Cục Map-Reduce}: Người dùng gửi câu hỏi toàn cục: Câu hỏi nhân bản thành 4 tia sáng chiếu đồng thời vào 4 báo cáo cộng đồng (Pha Map), từng cụm trả về điểm hữu ích và luận điểm, sau đó hội tụ về nút trung tâm để tổng hợp thành câu trả lời hoàn chỉnh (Pha Reduce).
  \item Render trực tiếp với lệnh: \code{manim -pql scene.py GraphRAGGlobalSearchScene}.''
\end{enumerate}
\end{quote}
\end{simulationbox}

---

\section{Hệ Thống Bài Tập \& Câu Hỏi Tư Duy Đa Tầng}
\label{sec:ch08-exercises}

\subsection{Nhóm A: Câu Hỏi Tư Duy Bản Chất (Conceptual Reasoning)}
\begin{enumerate}
  \item \textbf{Câu A.1}: Tại sao việc chuẩn hóa độ dài véc-tơ ($L_2$ Normalization) trước khi lưu vào Vector Database lại biến phép tính Cosine Similarity thành phép tính Dot Product? Lợi ích phần cứng của việc này trên CPU máy chủ là gì?
  \item \textbf{Câu A.2}: Phân tích nguyên nhân tại sao thuật toán HNSW bắt buộc phải duy trì toàn bộ cấu trúc chỉ mục trên RAM vật lý, và tại sao việc lưu chỉ mục HNSW trên ổ đĩa SSD truyền thống lại dẫn đến sụp đổ hiệu năng truy vấn.
  \item \textbf{Câu A.3}: Trong kỹ thuật Parent-Document Retrieval, tại sao chúng ta không nhúng trực tiếp cả đoạn mẹ và đoạn con vào cùng một không gian véc-tơ để tìm kiếm, mà lại tách biệt rõ ràng giữa chỉ mục tìm kiếm và dữ liệu trả về?
  \item \textbf{Câu A.4 (GraphRAG vs Vector RAG)}: Tại sao Vector RAG thuần túy lại thất bại trước các câu hỏi dạng suy luận toàn cục (Global Sensemaking)? GraphRAG sử dụng cơ chế nào để tổng hợp thông tin ở tầm vĩ mô?
  \item \textbf{Câu A.5 (Dense vs Sparse vs ColBERT Late Interaction)}: So sánh sự đánh đổi về độ trễ suy luận, chi phí lưu trữ bộ nhớ và độ chính xác giữa: Bi-Encoder (Dense), BM25 (Sparse) và ColBERT (Late Interaction).
  \item \textbf{Câu A.6 (Bản Chất Adaptive RAG)}: Tại sao việc bổ sung bước Retrieval Evaluator trong Corrective RAG (CRAG) lại ngăn chặn được hiện tượng ảo giác dây chuyền (\textit{Cascade Hallucination})?
  \item \textbf{Câu A.7 (Tại Sao RRF Ưu Việt Hơn Linear Score Fusion)}: \\
  Tại sao công thức cộng điểm tuyến tính $S = \alpha S_{\text{Dense}} + (1-\alpha) S_{\text{BM25}}$ thường thất bại trong môi trường sản xuất do hiện tượng trôi dạt phân phối điểm số (Score Distribution Drift)? Thuật toán Reciprocal Rank Fusion (RRF) giải quyết triệt để khuyết tật này như thế nào mà không cần chuẩn hóa điểm số?
  \item \textbf{Câu A.8 (Cơ Chế HyDE \& Điểm Mù Khi Triển Khai)}: \\
  Giải thích cơ chế hoạt động của kỹ thuật Hypothetical Document Embeddings (HyDE) trong việc khắc phục hiện tượng Bất đối xứng Ngữ nghĩa (Query-Document Asymmetry). Trong những tình huống nghiệp vụ nào (ví dụ: tra cứu văn bản pháp lý mới ban hành, mã lỗi hệ thống nội bộ) HyDE lại phản tác dụng và gây hại nghiêm trọng cho hệ thống RAG?
  \item \textbf{Câu A.9 (So Sánh Sâu Bi-Encoder, Cross-Encoder \& ColBERT)}: \\
  Phân tích tại sao Bi-Encoder gặp phải hiện tượng thắt cổ chai thông tin (\term{Information Bottleneck}) khi nén tài liệu vào một véc-tơ duy nhất. ColBERT giải quyết vấn đề này như thế nào mà vẫn đảm bảo khả năng tính toán trước ngoại tuyến (Offline Indexing)?
  \item \textbf{Câu A.10 (Contextual Retrieval \& Prompt Caching)}: \\
  Tại sao hiện tượng cô lập đoạn văn bản (\term{Chunk Isolation}) khiến các hệ thống Naive RAG thất bại khi tìm kiếm các dữ kiện cụ thể? Trình bày cách Contextual Retrieval khắc phục nhược điểm này và vai trò then chốt của Prompt Caching trong việc giảm chi phí vận hành.
\end{enumerate}

\subsection{Nhóm B: Bài Toán Tính Toán Giải Tích Bằng Tay (Analytical Problems)}
\begin{enumerate}
  \item \textbf{Bài B.1 (Tính Định Cỡ Bộ Nhớ RAM Cho Chỉ Mục HNSW)}:
  Doanh nghiệp cần xây dựng hệ thống tìm kiếm cho $N = 10.000.000$ (10 triệu) văn bản.
  Mỗi văn bản được nhúng thành một véc-tơ dạng FP32 gồm $d = 1.536$ chiều (chuẩn OpenAI text-embedding-3-large).
  Chỉ mục HNSW được cấu hình với số liên kết tối đa $M = 32$.
  Giả sử mỗi liên kết cạnh của đồ thị chiếm $8$ bytes con trỏ địa chỉ bộ nhớ 64-bit.
  \begin{itemize}
    \item[a)] Tính dung lượng RAM cần thiết để lưu trữ thuần túy các véc-tơ dữ liệu gốc.
    \item[b)] Tính dung lượng RAM xấp xỉ chỉ riêng cho cấu trúc đồ thị liên kết HNSW (bỏ qua các tầng trên, chỉ tính tầng Layer 0 nơi mỗi node có trung bình $M$ liên kết).
    \item[c)] Tính tổng dung lượng RAM tối thiểu của máy chủ để nạp toàn bộ cơ sở dữ liệu trên.
  \end{itemize}

  \item \textbf{Bài B.2 (Chứng Minh Tính Chất Hội Tụ Của RRF)}:
  Xét công thức RRF với hằng số làm mịn $k$:
  \begin{equation*}
  \text{RRF\_Score}(d) = \sum_{m=1}^M \frac{1}{k + r_m(d)}
  \end{equation*}
  Chứng minh rằng khi $k \to \infty$, thứ tự xếp hạng của RRF tiệm cận chính xác thứ tự xếp hạng của thuật toán tính trung bình thứ hạng (Mean Rank Borda Count): $\text{Score}_{\text{Borda}}(d) = -\sum_{m=1}^M r_m(d)$.

  \item \textbf{Bài B.3 (Tính Toán Chỉ Số Đánh Giá Ragas Bằng Tay)}:
  Một câu hỏi của người dùng nhận được câu trả lời từ LLM gồm 4 khẳng định (\textit{Atomic Claims}):
  \begin{enumerate}
    \item $C_1$: Doanh thu quý 3 của công ty đạt 150 tỷ đồng.
    \item $C_2$: Lợi nhuận ròng tăng trưởng $25\%$ so với cùng kỳ năm trước.
    \item $C_3$: Ban giám đốc quyết định chia cổ tức bằng tiền mặt tỷ lệ $10\%$.
    \item $C_4$: Công ty mở rộng thêm 5 chi nhánh tại khu vực miền Trung.
  \end{enumerate}
  Đối chiếu với ngữ cảnh trích xuất từ 3 tài liệu, chuyên viên xác định $C_1$, $C_2$, $C_4$ đều có bằng chứng trực tiếp trong tài liệu, riêng $C_3$ hoàn toàn không được đề cập trong tài liệu.
  Hãy tính chỉ số \textbf{Faithfulness} của câu trả lời trên. Nếu người dùng yêu cầu chỉ số Faithfulness tối thiểu phải đạt $0.90$ để phê duyệt xuất bản tự động, hệ thống có được phép trả lời người dùng hay không?

  \item \textbf{Bài B.4 (Tính Toán Thủ Công Điểm Số BM25 Chuẩn Okapi)}: \\
  Cho tập dữ liệu gồm $N = 3$ tài liệu với độ dài lần lượt là $|D_1| = 4$ từ, $|D_2| = 6$ từ, $|D_3| = 5$ từ. Độ dài trung bình $\text{avgdl} = 5.0$ từ.
  Từ khóa truy vấn $q = \{\text{``AI''}\}$.
  Từ ``AI'' xuất hiện trong $D_1$ với tần số $f_1 = 2$, trong $D_2$ với tần số $f_2 = 1$, và không xuất hiện trong $D_3$ ($n(q) = 2$).
  Tham số chuẩn: $k_1 = 1.2, b = 0.75$.
  \begin{enumerate}
    \item Tính giá trị $\text{IDF}(\text{``AI''})$ theo công thức Okapi BM25 chuẩn:
    \begin{equation*}
    \text{IDF}(q) = \ln \left( \frac{N - n(q) + 0.5}{n(q) + 0.5} + 1 \right)
    \end{equation*}
    \item Tính điểm số $\text{BM25}(D_1, q)$ và $\text{BM25}(D_2, q)$. Xác định tài liệu nào có độ phù hợp cao hơn.
  \end{enumerate}

  \item \textbf{Bài B.5 (Tính Toán Điểm Số Reciprocal Rank Fusion - RRF)}: \\
  Cho 4 tài liệu $\{d_1, d_2, d_3, d_4\}$ được truy xuất từ hai hệ thống:
  \begin{itemize}
    \item Hệ thống Dense Vector Search xếp hạng: $r_{\text{dense}}(d_1) = 1$, $r_{\text{dense}}(d_2) = 2$, $r_{\text{dense}}(d_3) = 3$, $r_{\text{dense}}(d_4) = 4$.
    \item Hệ thống BM25 Sparse Search xếp hạng: $r_{\text{bm25}}(d_3) = 1$, $r_{\text{bm25}}(d_1) = 2$, $r_{\text{bm25}}(d_4) = 3$, $r_{\text{bm25}}(d_2) = 4$.
  \end{itemize}
  Với hằng số làm mịn $k = 60$:
  \begin{enumerate}
    \item Tính toán điểm số $\text{RRF}(d_i)$ cho cả 4 tài liệu bằng tay.
    \item Xác định thứ tự xếp hạng dung hợp cuối cùng và cho biết tài liệu nào chiến thắng.
  \end{enumerate}

  \item \textbf{Bài B.6 (Tính Toán Thủ Công Ma Trận Tương Tác Muộn MaxSim ColBERT)}: \\
  Cho câu truy vấn $q$ gồm 2 tokens với ma trận biểu diễn $E_q = \begin{bmatrix} 1 & 0 & 0 & 0 \\ 0 & 1 & 0 & 0 \end{bmatrix}$ (kích thước $2 \times 4$).
  Tài liệu $D$ gồm 3 tokens với ma trận biểu diễn:
  \begin{equation*}
  E_d = \begin{bmatrix} 0 & 1 & 0 & 0 \\ 1 & 0 & 0 & 0 \\ 0 & 0 & 1 & 0 \end{bmatrix} \in \mathbb{R}^{3 \times 4}
  \end{equation*}
  Tất cả các véc-tơ đều đã được chuẩn hóa độ dài Euclid $\|\mathbf{e}\|_2 = 1$.
  \begin{itemize}
    \item[a)] Tính ma trận tương đồng tích vô hướng $S = E_q E_d^T \in \mathbb{R}^{2 \times 3}$.
    \item[b)] Xác định giá trị cực đại $\max_{j} S_{i, j}$ cho từng token truy vấn $i \in \{0, 1\}$.
    \item[c)] Tính điểm số $\text{MaxSim}(q, D)$ và giải thích ý nghĩa so khớp ngữ nghĩa mềm (Soft Alignment).
  \end{itemize}

  \item \textbf{Bài B.7 (Phân Tích Chi Phí \& Độ Phức Tạp Map-Reduce Global Search Trong GraphRAG)}: \\
  Xét một kho tài liệu $5.000$ trang được chia thành $20.000$ chunks.
  Khi xây dựng đồ thị tri thức, thuật toán Leiden phát hiện được $K = 50$ cộng đồng ở Cấp độ 0 (C0 Macro) và $M = 300$ cộng đồng ở Cấp độ 1 (C1 Meso).
  Mỗi bản tóm tắt cộng đồng C0 có độ dài trung bình $1.000$ tokens, câu hỏi của người dùng dài $50$ tokens.
  \begin{itemize}
    \item[a)] Khi thực hiện truy vấn toàn cục ở Cấp độ 0, pha Map cần gửi bao nhiêu request song song tới LLM và tổng số tokens đầu vào của toàn bộ pha Map là bao nhiêu?
    \item[b)] Nếu mỗi lượt Map trích xuất $5$ luận điểm cốt lõi (mỗi luận điểm dài $30$ tokens) kèm điểm số hữu ích, và pha Reduce chọn lọc top $20$ luận điểm có điểm cao nhất để tổng hợp câu trả lời, hãy tính kích thước ngữ cảnh (tokens) mà pha Reduce đưa vào LLM. So sánh tổng chi phí token này với việc nhồi nhét trực tiếp $20.000$ chunks vào một cửa sổ ngữ cảnh khổng lồ.
  \end{itemize}
\end{enumerate}

\subsection{Nhóm C: Thực Hành Kỹ Nghệ \& Lập Trình (Implementation Labs)}
\begin{enumerate}
  \item \textbf{Bài C.1 (Bộ Tái Xếp Hạng Hai Tầng Hoàn Chỉnh)}:
  Hiện thực hóa lớp \code{TwoStageHybridRetriever} kết hợp BM25 và Dense Search qua RRF, sau đó dùng Cross-Encoder để tái xếp hạng Top-5 kết quả cuối cùng. Viết unit test chứng minh tài liệu đứng hạng 3 ở Stage 1 có thể leo lên hạng 1 ở Stage 2.

  \item \textbf{Bài C.2 (Bộ Đánh Giá Độ Trung Thực Faithfulness)}:
  Hiện thực hóa lớp \code{FaithfulnessEvaluator} nhận vào danh sách các câu khẳng định và văn bản ngữ cảnh, tự động tính toán chỉ số Faithfulness theo tỷ lệ các khẳng định có mặt trong ngữ cảnh.

  \item \textbf{Bài C.3 (Bộ Truy Xuất Tri Thức Đồ Thị k-hop)}:
  Xây dựng lớp \code{GraphEntityRetriever} lưu danh sách kề của đồ thị tri thức doanh nghiệp. Viết hàm \code{expand_subgraph(seed_entities, k_hops)} mở rộng không gian ngữ cảnh qua $k$ bước nhảy đồ thị để phục vụ bổ sung ngữ cảnh cho LLM.

  \item \textbf{Bài C.4 (Bộ Chia Đoạn Bảng Biểu Nhận Biết Cấu Trúc - TableAwareChunker)}: \\
  Xây dựng lớp Python \code{TableAwareChunker} có khả năng tự động phát hiện cấu trúc bảng biểu Markdown (\code{| col1 | col2 |}). Khi bảng biểu có số dòng vượt quá giới hạn $N$ dòng, thuật toán tự động tách bảng thành các chunk con nhưng luôn sao chép dòng tiêu đề cột vào đầu mỗi chunk con để bảo toàn trọn vẹn ngữ nghĩa số liệu cho mô hình embedding.

  \item \textbf{Bài C.5 (Hiện Thực Hóa Động Cơ ColBERT MaxSim)}: \\
  Xây dựng lớp \code{ColBERTMaxSimEngine} với phương thức tĩnh \code{compute_maxsim} nhận vào hai ma trận NumPy đã chuẩn hóa $L_2$ và tính toán điểm tương đồng MaxSim. Viết unit test xác thực kết quả số học khớp chính xác với Bài B.6.

  \item \textbf{Bài C.6 (Bộ Tiền Xử Lý Ngữ Cảnh Hóa ContextualChunkAugmenter)}: \\
  Xây dựng lớp \code{ContextualChunkAugmenter} nhận từ điển ngữ cảnh tài liệu và chèn cấu trúc thẻ XML \code{<document_context>} vào trước từng chunk thô, chuẩn bị dữ liệu đầu vào tối ưu cho cả BM25 và Vector Database.
\end{enumerate}

\subsection{Nhóm D: Phân Tích Thất Bại \& Debugging Thực Tế (Post-Mortem Analysis)}
\begin{enumerate}
  \item \textbf{Sự Cố 1 (Xung Đột Phiên Bản Tài Liệu / Lỗi Trôi Thời Gian)}:
  Trợ lý RAG của ngân hàng trả về lãi suất cũ $7.2\%$ thay vì thông tư mới $5.5\%$, mặc dù cả hai văn bản đều có trong Vector DB và văn bản cũ có điểm Cosine cao hơn. Hãy phân tích nguyên nhân gốc rễ và đề xuất giải pháp khắc phục triệt để.

  \item \textbf{Sự Cố 2 (Lỗi Chunking Cắt Ngang Bảng Biểu Tài Chính)}:
  Một công ty chứng khoán triển khai RAG phân tích báo cáo tài chính PDF. Người dùng hỏi: ``Doanh thu mảng bán lẻ của công ty năm 2025 là bao nhiêu?''. Hệ thống trả về con số của năm 2024!
  Kiểm tra log cho thấy bộ chia nhỏ văn bản (\textit{Text Splitter}) dùng thuật toán cắt cố định $512$ ký tự đã vô tình cắt ngang bảng biểu Markdown: Hàng tiêu đề của bảng (\code{| Năm 2024 | Năm 2025 |}) nằm ở Chunk 1, trong khi các dòng số liệu (\code{| 1.200 tỷ | 1.850 tỷ |}) lại nằm ở Chunk 2! Khi Vector Search trích xuất Chunk 2, LLM hoàn toàn mất tiêu đề cột và tự gán sai số liệu!
  Hãy viết báo cáo sự cố (Post-Mortem) và thiết kế giải pháp phân đoạn văn bản nhận biết cấu trúc (\textit{Table-Aware Chunking}).
\end{enumerate}

---

\section{Hướng Dẫn Lời Giải \& Bình Luận Kỹ Thuật Chi Tiết}
\label{sec:ch08-solutions}

\subsection{Lời Giải Nhóm A: Câu Hỏi Tư Duy Bản Chất}
\begin{itemize}
  \item \textbf{Lời giải Câu A.1}:
  Khoảng cách Cosine giữa $\mathbf{u}$ và $\mathbf{v}$ là $1 - \frac{\mathbf{u} \cdot \mathbf{v}}{\|\mathbf{u}\|_2 \|\mathbf{v}\|_2}$. Khi cả hai véc-tơ đã được chuẩn hóa độ dài $\|\mathbf{u}\|_2 = \|\mathbf{v}\|_2 = 1$, mẫu số bằng $1$, và phép đo tương đồng Cosine thu gọn chính xác về tích vô hướng $\mathbf{u} \cdot \mathbf{v}$.
  Trên phần cứng CPU, việc tính căn bậc hai và phép chia của chuẩn Euclid là những chỉ lệnh tốn rất nhiều chu kỳ xung nhịp (Clock Cycles). Ngược lại, tích vô hướng là phép nhân-cộng dồn (Multiply-Accumulate -- FMA) được tăng tốc hoàn hảo bởi tập lệnh SIMD (AVX-512), cho phép xử lý song song 16 phép tính số thực trong đúng 1 chu kỳ xung nhịp, tăng tốc độ truy vấn lên từ $5\times$ đến $10\times$ lần.

  \item \textbf{Lời giải Câu A.2}:
  Thuật toán HNSW hoạt động bằng cách nhảy liên tục qua các đỉnh của đồ thị đa chiều. Mỗi bước nhảy đòi hỏi việc đọc ngẫu nhiên (Random Memory Access) danh sách các node láng giềng liên kết với node hiện tại.
  Thời gian truy xuất ngẫu nhiên trên RAM chỉ mất khoảng $50\text{ đến }100\text{ nano-giây}$. Ngược lại, thời gian truy xuất ngẫu nhiên trên ổ SSD NVMe nhanh nhất vẫn mất từ $10\text{ đến }20\text{ micro-giây}$ (chậm hơn $100\text{ đến }200$ lần!). Một lượt tìm kiếm HNSW cần nhảy qua vài trăm đỉnh đồ thị; nếu dữ liệu nằm trên SSD, độ trễ sẽ bị kéo dài từ vài mili-giây lên hàng trăm mili-giây, làm tê liệt hoàn toàn hệ thống thời gian thực.

  \item \textbf{Lời giải Câu A.3}:
  Nếu nhúng cả đoạn mẹ ($1.024$ tokens) và đoạn con ($128$ tokens) vào cùng một không gian véc-tơ, sẽ xảy ra sự cạnh tranh bất bình đẳng. Đoạn mẹ chứa nhiều ý tưởng hỗn hợp sẽ có véc-tơ bị kéo về phía trung tâm của không gian, trong khi đoạn con có véc-tơ tập trung ở các vùng rìa sắc nhọn. Việc này làm phá vỡ phân phối khoảng cách. Bằng cách chỉ nhúng các đoạn con, ta duy trì được một không gian metric đồng nhất và sắc nét nhất cho thuật toán tìm kiếm, trong khi việc dùng con trỏ Metadata lấy đoạn mẹ đảm bảo cung cấp ngữ cảnh đầy đủ nhất cho LLM.

  \item \textbf{Lời giải Câu A.4}:
  Vector RAG tìm kiếm dựa trên độ tương đồng giữa câu hỏi và từng chunk văn bản độc lập. Khi câu hỏi mang tính tổng hợp vĩ mô, không có chunk nào chứa trọn vẹn câu trả lời. GraphRAG giải quyết điểm mù này bằng cách trích xuất đồ thị thực thể-quan hệ, sau đó dùng thuật toán Leiden phân cụm đồ thị và sinh các báo cáo tóm tắt phân cấp (Hierarchical Community Summaries). Truy vấn toàn cục sẽ map-reduce các bản tóm tắt cộng đồng vĩ mô này, mang lại câu trả lời bao quát toàn diện.

  \item \textbf{Lời giải Câu A.5}:
  - \textbf{Bi-Encoder (Dense)}: Nén toàn bộ đoạn văn bản thành 1 véc-tơ duy nhất. Tốc độ tìm kiếm ANN cực nhanh ($O(\log N)$ với HNSW), bộ nhớ vừa phải, nhưng mất mát thông tin cục bộ và kém chính xác với từ khóa hiếm.
  - \textbf{Sparse (BM25)}: Dựa trên tần số từ khóa đảo (Inverted Index). Không cần GPU, tốc độ cực nhanh, cực chuẩn với mã số và từ khóa chính xác, nhưng mù tịt về ngữ nghĩa đồng nghĩa.
  - \textbf{ColBERT (Late Interaction)}: Giữ lại véc-tơ nhúng của từng token riêng biệt và tính tương đồng MaxSim. Độ chính xác tiệm cận Cross-Encoder nhưng dung lượng lưu trữ bộ nhớ tăng gấp hàng chục lần so với Bi-Encoder vì phải lưu hàng triệu token embeddings.

  \item \textbf{Lời giải Câu A.6}:
  Trong Naive RAG, nếu bước truy xuất trả về tài liệu rác, LLM vẫn buộc phải đọc và cố gắng sinh câu trả lời, dẫn tới ảo giác (Hallucination). Corrective RAG (CRAG) chèn thêm một bộ đánh giá chất lượng (Retrieval Evaluator). Nếu điểm tin cậy thấp, hệ thống ngắt luồng sinh câu trả lời từ tài liệu rác và tự động chuyển hướng (Fallback) sang Web Search hoặc đưa ra thông báo từ chối an toàn, dập tắt ảo giác ngay từ đầu vào.

  \item \textbf{Lời giải Câu A.7}:
  Phương pháp cộng điểm tuyến tính $S = \alpha S_{\text{Dense}} + (1-\alpha) S_{\text{BM25}}$ đòi hỏi hai hệ thống phải có cùng thang đo điểm số. Tuy nhiên:
  - Điểm BM25 không bị chặn trên ($[0, +\infty)$) và phụ thuộc mạnh vào độ dài câu truy vấn cũng như tần số từ khóa trong tài liệu.
  - Điểm Dense Cosine nằm trong khoảng $[-1, 1]$ (hoặc $[0, 1]$ sau khi chuẩn hóa).
  Nếu chỉ dùng Min-Max Normalization trên tập kết quả Top-K cục bộ, một tài liệu ngoại lai có điểm BM25 cực cao sẽ nén phẳng toàn bộ điểm của các tài liệu khác về gần 0.
  Reciprocal Rank Fusion (RRF) giải quyết triệt để vấn đề này vì nó \textbf{hoàn toàn loại bỏ giá trị điểm số tuyệt đối} và chỉ thao tác trên \textbf{thứ hạng tương đối} $r_m(d) \in \{1, 2, \dots, K\}$. Nhờ hằng số làm mịn $k \approx 60$, RRF đóng vai trò như một bộ lọc thông thấp (Low-pass filter), làm giảm độ nhạy cảm của các vị trí đầu bảng và dung hợp công bằng kết quả giữa các mô hình tìm kiếm có phân phối điểm số hoàn toàn khác biệt.

  \item \textbf{Lời giải Câu A.8}:
  HyDE (Hypothetical Document Embeddings) khắc phục hiện tượng Bất đối xứng Ngữ nghĩa bằng cách yêu cầu LLM sinh ra một đoạn văn bản giả định trả lời câu hỏi trước, sau đó nhúng đoạn văn bản giả định này vào không gian véc-tơ để tìm kiếm các tài liệu thực tế tương tự. Vì cả hai đều là văn bản dài có cấu trúc tương đương, khoảng cách ngữ nghĩa được thu hẹp đáng kể.
  Tuy nhiên, HyDE sẽ phản tác dụng nghiêm trọng trong hai tình huống:
  \begin{itemize}
    \item \textbf{Tri thức mới hoặc bí mật nội bộ}: Khi người dùng hỏi về quy chế mới ban hành hôm qua hoặc mã lỗi nội bộ chưa từng xuất hiện trong dữ liệu huấn luyện của LLM, mô hình sẽ sinh ra văn bản giả định chứa đầy ảo giác sai lệch. Việc nhúng văn bản ảo giác này sẽ dẫn đường cho Vector Search trích xuất các tài liệu hoàn toàn sai!
    \item \textbf{Gia tăng độ trễ và chi phí}: Mỗi truy vấn phải chịu thêm độ trễ TTFT của một lượt sinh văn bản hoàn chỉnh từ LLM (thường mất từ $500\text{ ms}$ đến $2\text{ giây}$), phá hủy yêu cầu phản hồi thời gian thực của ứng dụng tìm kiếm.
  \end{itemize}

  \item \textbf{Lời giải Câu A.9}:
  \textbf{Bi-Encoder} nén toàn bộ chuỗi token của tài liệu thành một véc-tơ duy nhất $\mathbf{e}_d \in \mathbb{R}^d$ thông qua cơ chế Pooling (Mean Pooling hoặc CLS Token). Thao tác này tạo ra nút thắt cổ chai thông tin (\term{Information Bottleneck}): mọi chi tiết cú pháp cục bộ, số liệu cụ thể và từ khóa chuyên ngành hiếm đều bị hòa tan vào một điểm duy nhất trong không gian đa chiều, làm mất khả năng phân biệt giữa các sắc thái ngữ nghĩa tinh vi.
  \textbf{ColBERT} giải quyết triệt để vấn đề này bằng mô hình Tương Tác Muộn (\term{Late Interaction}): Thay vì gộp véc-tơ, nó duy trì ma trận embedding của từng token riêng biệt $E_d \in \mathbb{R}^{|d| \times d}$. Điểm mấu chốt là $E_d$ hoàn toàn độc lập với câu truy vấn, cho phép tính toán trước ngoại tuyến $100\%$ và lưu trữ vào chỉ mục. Khi có truy vấn $q$, hệ thống tính $E_q$ và áp dụng toán tử MaxSim $\sum_{i} \max_{j} (E_{q, i} \cdot E_{d, j}^T)$. Nhờ đó, ColBERT bảo toàn trọn vẹn độ chính xác tương tác cấp token tiệm cận Cross-Encoder mà vẫn duy trì tốc độ truy vấn mili-giây như Bi-Encoder.

  \item \textbf{Lời giải Câu A.10}:
  Hiện tượng \term{Chunk Isolation} (Cô lập đoạn văn bản) xảy ra khi tài liệu dài bị cắt nhỏ thành các chunk $300 - 500$ tokens độc lập. Nhiều chunk nằm giữa tài liệu chứa các đại từ thay thế (như ``công ty này'', ``nó'', ``trong năm vừa qua'') hoặc bảng số liệu mà không có tên thực thể hay mốc thời gian cụ thể. Cả BM25 lẫn Dense Search đều thất bại khi người dùng tìm kiếm từ khóa cụ thể.
  \textbf{Contextual Retrieval} (Anthropic, 2024) khắc phục bằng cách dùng LLM tạo một đoạn tóm tắt ngữ cảnh $50 - 100$ tokens gắn vào đầu mỗi chunk trước khi lập chỉ mục. Đòn bẩy \textbf{Prompt Caching} đóng vai trò then chốt vì toàn bộ tài liệu gốc được lưu vào KV Cache tĩnh trên GPU; các lệnh sinh ngữ cảnh cho từng chunk chỉ cần đọc lại KV Cache với chi phí giảm tới $90\%$, biến giải pháp này thành khả thi về mặt kinh tế cho quy mô triệu văn bản.
\end{itemize}

\subsection{Lời Giải Nhóm B: Bài Toán Tính Toán Giải Tích Bằng Tay}

\paragraph{Lời giải Bài B.1}:
\begin{itemize}
  \item[a)] Dung lượng RAM chứa véc-tơ dữ liệu gốc:
  Mỗi véc-tơ $1.536$ chiều dạng FP32 tốn: $1.536 \times 4\text{ Bytes} = 6.144\text{ Bytes} = 6\text{ KiB}$.
  Với 10 triệu véc-tơ:
  \begin{equation*}
  V_{\text{raw}} = 10^7 \times 6.144\text{ Bytes} = 61.440.000.000\text{ Bytes} \approx 57.22\text{ GiB}
  \end{equation*}

  \item[b)] Dung lượng RAM cho các cạnh liên kết của đồ thị HNSW:
  Tại Layer 0, mỗi node duy trì trung bình $M = 32$ liên kết cạnh tới các node khác.
  Mỗi liên kết cạnh là một con trỏ địa chỉ 64-bit ($8$ bytes).
  Bộ nhớ liên kết cho 1 node: $32 \times 8\text{ Bytes} = 256\text{ Bytes}$.
  Với 10 triệu node:
  \begin{equation*}
  V_{\text{graph}} = 10^7 \times 256\text{ Bytes} = 2.560.000.000\text{ Bytes} \approx 2.38\text{ GiB}
  \end{equation*}
  (Nếu tính thêm cấu trúc bảng băm, metadata ID, và các tầng trên, chi phí này trong thực tế thường dao động từ $4\text{ đến }8\text{ GiB}$).

  \item[c)] Tổng dung lượng RAM tối thiểu:
  \begin{equation*}
  V_{\text{total}} = V_{\text{raw}} + V_{\text{graph}} + V_{\text{OS/Buffer}} \approx 57.22\text{ GiB} + 2.38\text{ GiB} + 12\text{ GiB} \approx 72\text{ GiB}
  \end{equation*}
  \textit{Khuyến nghị kỹ thuật}: Doanh nghiệp cần trang bị máy chủ có ít nhất \textbf{$128\text{ GB}$ RAM ECC} để đảm bảo hệ thống hoạt động ổn định và có bộ đệm hấp thụ tải.
\end{itemize}

\paragraph{Lời giải Bài B.2}:
Khai triển chuỗi Taylor của số hạng RRF khi $k \to \infty$:
\begin{equation*}
\frac{1}{k + r_m(d)} = \frac{1}{k \left(1 + \frac{r_m(d)}{k}\right)} = \frac{1}{k} \left( 1 - \frac{r_m(d)}{k} + O\left(\frac{1}{k^2}\right) \right) = \frac{1}{k} - \frac{r_m(d)}{k^2} + O\left(\frac{1}{k^3}\right)
\end{equation*}
Lấy tổng trên toàn bộ $M$ hệ thống xếp hạng:
\begin{equation*}
\text{RRF\_Score}(d) = \sum_{m=1}^M \left( \frac{1}{k} - \frac{r_m(d)}{k^2} \right) + O\left(\frac{1}{k^3}\right) = \frac{M}{k} - \frac{1}{k^2} \sum_{m=1}^M r_m(d) + O\left(\frac{1}{k^3}\right)
\end{equation*}
Vì $\frac{M}{k}$ là hằng số giống nhau cho mọi tài liệu $d$, và hệ số $\frac{1}{k^2} > 0$, việc tối đa hóa $\text{RRF\_Score}(d)$ tương đương với việc tối thiểu hóa tổng thứ hạng $\sum_{m=1}^M r_m(d)$ (tức tối đa hóa điểm Borda Count âm $-\sum r_m(d)$). (Điều phải chứng minh).

\paragraph{Lời giải Bài B.3}:
Tổng số khẳng định có trong câu trả lời: $|\text{Total Claims}| = 4$.
Số khẳng định có bằng chứng trực tiếp trong ngữ cảnh trích xuất: $|\text{Supported Claims}| = 3$ ($C_1, C_2, C_4$).
Chỉ số Faithfulness:
\begin{equation*}
\text{Faithfulness} = \frac{3}{4} = 0.75 \quad (\mathbf{75\%})
\end{equation*}
Vì $0.75 < 0.90$ (ngưỡng tối thiểu của doanh nghiệp), hệ thống \textbf{tuyệt đối không được phép tự động xuất bản câu trả lời này}! Hệ thống phải kích hoạt cờ cảnh báo hoặc yêu cầu chuyên viên thẩm định lại trước khi gửi cho khách hàng.

\paragraph{Lời giải Bài B.4}:
\begin{enumerate}
  \item Tính giá trị $\text{IDF}(\text{``AI''})$:
  Với $N = 3$, số tài liệu chứa từ khóa $n(q) = 2$:
  \begin{equation*}
  \text{IDF}(q) = \ln \left( \frac{3 - 2 + 0.5}{2 + 0.5} + 1 \right) = \ln \left( \frac{1.5}{2.5} + 1 \right) = \ln (0.6 + 1) = \ln(1.6) \approx \mathbf{0.4700}
  \end{equation*}

  \item Tính điểm số BM25 cho $D_1$ và $D_2$:
  Độ dài trung bình $\text{avgdl} = 5.0$.
  - Với $D_1$ ($|D_1| = 4, f_1 = 2$):
  \begin{equation*}
  B_1 = 1 - b + b \cdot \frac{|D_1|}{\text{avgdl}} = 1 - 0.75 + 0.75 \cdot \frac{4}{5.0} = 0.25 + 0.60 = 0.85
  \end{equation*}
  Số hạng tần số:
  \begin{equation*}
  \frac{f_1 (k_1 + 1)}{f_1 + k_1 \cdot B_1} = \frac{2 \times (1.2 + 1)}{2 + 1.2 \times 0.85} = \frac{4.4}{2 + 1.02} = \frac{4.4}{3.02} \approx 1.4570
  \end{equation*}
  Điểm số:
  \begin{equation*}
  \text{BM25}(D_1, q) = 0.4700 \times 1.4570 \approx \mathbf{0.6848}
  \end{equation*}

  - Với $D_2$ ($|D_2| = 6, f_2 = 1$):
  \begin{equation*}
  B_2 = 1 - 0.75 + 0.75 \cdot \frac{6}{5.0} = 0.25 + 0.90 = 1.15
  \end{equation*}
  Số hạng tần số:
  \begin{equation*}
  \frac{f_2 (k_1 + 1)}{f_2 + k_1 \cdot B_2} = \frac{1 \times 2.2}{1 + 1.2 \times 1.15} = \frac{2.2}{1 + 1.38} = \frac{2.2}{2.38} \approx 0.9244
  \end{equation*}
  Điểm số:
  \begin{equation*}
  \text{BM25}(D_2, q) = 0.4700 \times 0.9244 \approx \mathbf{0.4345}
  \end{equation*}
  \textbf{Kết luận}: $D_1$ có điểm BM25 cao hơn $D_2$ ($0.6848 > 0.4345$) vì tần số xuất hiện nhiều hơn và độ dài tài liệu ngắn hơn.
\end{enumerate}

\paragraph{Lời giải Bài B.5}:
Áp dụng công thức RRF với $k = 60$: $\text{RRF}(d) = \frac{1}{60 + r_{\text{dense}}(d)} + \frac{1}{60 + r_{\text{bm25}}(d)}$.
\begin{enumerate}
  \item Tính toán điểm số cho 4 tài liệu:
  \begin{itemize}
    \item Tài liệu $d_1$ ($r_{\text{dense}} = 1, r_{\text{bm25}} = 2$):
    \begin{equation*}
    \text{RRF}(d_1) = \frac{1}{60 + 1} + \frac{1}{60 + 2} = \frac{1}{61} + \frac{1}{62} \approx 0.016393 + 0.016129 = \mathbf{0.032522}
    \end{equation*}
    \item Tài liệu $d_2$ ($r_{\text{dense}} = 2, r_{\text{bm25}} = 4$):
    \begin{equation*}
    \text{RRF}(d_2) = \frac{1}{60 + 2} + \frac{1}{60 + 4} = \frac{1}{62} + \frac{1}{64} \approx 0.016129 + 0.015625 = \mathbf{0.031754}
    \end{equation*}
    \item Tài liệu $d_3$ ($r_{\text{dense}} = 3, r_{\text{bm25}} = 1$):
    \begin{equation*}
    \text{RRF}(d_3) = \frac{1}{60 + 3} + \frac{1}{60 + 1} = \frac{1}{63} + \frac{1}{61} \approx 0.015873 + 0.016393 = \mathbf{0.032266}
    \end{equation*}
    \item Tài liệu $d_4$ ($r_{\text{dense}} = 4, r_{\text{bm25}} = 3$):
    \begin{equation*}
    \text{RRF}(d_4) = \frac{1}{60 + 4} + \frac{1}{60 + 3} = \frac{1}{64} + \frac{1}{63} \approx 0.015625 + 0.015873 = \mathbf{0.031498}
    \end{equation*}
  \end{itemize}

  \item Thứ tự xếp hạng dung hợp RRF cuối cùng:
  \begin{equation*}
  \mathbf{d_1} \ (0.032522) > \mathbf{d_3} \ (0.032266) > \mathbf{d_2} \ (0.031754) > \mathbf{d_4} \ (0.031498)
  \end{equation*}
  Tài liệu \textbf{$d_1$ giành chiến thắng} vì duy trì thứ hạng rất cao ở cả hai hệ thống ($1$ và $2$), vượt qua $d_3$ vốn đứng đầu ở BM25 nhưng tụt xuống hạng $3$ ở Dense Search!
\end{enumerate}

\paragraph{Lời giải Bài B.6}:
\begin{itemize}
  \item[a)] Tính ma trận tương đồng $S = E_q E_d^T$:
  Ta có ma trận $E_d^T \in \mathbb{R}^{4 \times 3}$:
  \begin{equation*}
  E_d^T = \begin{bmatrix} 0 & 1 & 0 \\ 1 & 0 & 0 \\ 0 & 0 & 1 \\ 0 & 0 & 0 \end{bmatrix}
  \end{equation*}
  Thực hiện nhân ma trận $S = E_q E_d^T \in \mathbb{R}^{2 \times 3}$:
  \begin{itemize}
    \item Dòng 0: $E_{q, 0} = [1, 0, 0, 0]$:
    \begin{equation*}
    S_{0, 0} = 1 \cdot 0 + 0 \cdot 1 + 0 \cdot 0 + 0 \cdot 0 = 0
    \end{equation*}
    \begin{equation*}
    S_{0, 1} = 1 \cdot 1 + 0 \cdot 0 + 0 \cdot 0 + 0 \cdot 0 = 1
    \end{equation*}
    \begin{equation*}
    S_{0, 2} = 1 \cdot 0 + 0 \cdot 0 + 0 \cdot 1 + 0 \cdot 0 = 0
    \end{equation*}
    \item Dòng 1: $E_{q, 1} = [0, 1, 0, 0]$:
    \begin{equation*}
    S_{1, 0} = 0 \cdot 0 + 1 \cdot 1 + 0 \cdot 0 + 0 \cdot 0 = 1
    \end{equation*}
    \begin{equation*}
    S_{1, 1} = 0 \cdot 1 + 1 \cdot 0 + 0 \cdot 0 + 0 \cdot 0 = 0
    \end{equation*}
    \begin{equation*}
    S_{1, 2} = 0 \cdot 0 + 1 \cdot 0 + 0 \cdot 1 + 0 \cdot 0 = 0
    \end{equation*}
  \end{itemize}
  Kết quả ma trận tương đồng tích vô hướng:
  \begin{equation*}
  S = \begin{bmatrix} 0 & 1 & 0 \\ 1 & 0 & 0 \end{bmatrix} \in \mathbb{R}^{2 \times 3}
  \end{equation*}

  \item[b)] Xác định giá trị cực đại trên từng dòng:
  \begin{itemize}
    \item Token truy vấn $0$: $\max_{j \in \{0, 1, 2\}} S_{0, j} = \max(0, 1, 0) = 1.0$ (khớp tại token 1 của tài liệu).
    \item Token truy vấn $1$: $\max_{j \in \{0, 1, 2\}} S_{1, j} = \max(1, 0, 0) = 1.0$ (khớp tại token 0 của tài liệu).
  \end{itemize}

  \item[c)] Điểm số MaxSim:
  \begin{equation*}
  \text{MaxSim}(q, D) = \sum_{i=0}^1 \max_{j} S_{i, j} = 1.0 + 1.0 = \mathbf{2.0}
  \end{equation*}
  \textit{Ý nghĩa}: Cả 2 token của câu truy vấn đều tìm thấy token đối ứng chính xác tuyệt đối trong tài liệu (dù thứ tự bị đảo ngược: token $0$ của query khớp token $1$ của doc, token $1$ của query khớp token $0$ của doc). ColBERT cho phép khớp ngữ nghĩa mềm mà không đòi hỏi cứng nhắc về thứ tự cú pháp hay nén phẳng thông tin.
\end{itemize}

\paragraph{Lời giải Bài B.7}:
\begin{itemize}
  \item[a)] Số request và tokens trong pha Map ở Cấp độ 0 (C0 Macro):
  Vì có $K = 50$ cộng đồng C0, pha Map cần gửi đúng \textbf{$50$ requests độc lập song song} tới LLM.
  Mỗi request gồm câu hỏi ($50$ tokens) và bản tóm tắt cộng đồng ($1.000$ tokens):
  \begin{equation*}
  L_{\text{request}} = 50 + 1.000 = 1.050\text{ tokens}
  \end{equation*}
  Tổng số tokens đầu vào của toàn bộ pha Map:
  \begin{equation*}
  N_{\text{map\_in}} = 50 \times 1.050 = \mathbf{52.500\text{ tokens}}
  \end{equation*}

  \item[b)] Kích thước ngữ cảnh pha Reduce và so sánh:
  Mỗi lượt Map trích xuất $5$ luận điểm, tổng cộng có $50 \times 5 = 250$ luận điểm.
  Pha Reduce chỉ gom top $20$ luận điểm có điểm số hữu ích cao nhất:
  \begin{equation*}
  L_{\text{key\_points}} = 20 \times 30 = 600\text{ tokens}
  \end{equation*}
  Cộng thêm câu hỏi gốc ($50$ tokens) và câu lệnh điều phối ($150$ tokens), kích thước ngữ cảnh nạp vào LLM ở pha Reduce chỉ vỏn vẹn:
  \begin{equation*}
  L_{\text{reduce\_prompt}} = 600 + 50 + 150 = \mathbf{800\text{ tokens}}
  \end{equation*}
  \textit{So sánh định lượng}:
  Nếu nhồi nhét trực tiếp $20.000$ chunks (mỗi chunk $300$ tokens) vào một cửa sổ ngữ cảnh khổng lồ, ta cần:
  \begin{equation*}
  N_{\text{naive}} = 20.000 \times 300 = \mathbf{6.000.000\text{ tokens (6 triệu tokens!)}}
  \end{equation*}
  GraphRAG đã nén lượng thông tin cần đọc từ 6 triệu tokens xuống chỉ còn $52.500$ tokens ở pha Map và $800$ tokens ở pha Reduce (giảm tới \textbf{$114\times$ lần}), đồng thời loại bỏ hoàn toàn hiện tượng ``Lạc lối ở giữa'' (Lost in the middle) để sinh ra câu trả lời toàn cục chất lượng cao.
\end{itemize}

\subsection{Lời Giải Nhóm C: Thực Hành Kỹ Nghệ \& Lập Trình}
Toàn bộ mã nguồn kiểm chứng đã được cung cấp tại Mục~\ref{sec:python-hybrid-engine} với đầy đủ các module:
\begin{enumerate}
  \item \code{TwoStageHybridRetriever}: Kết hợp BM25 và Dense Search qua RRF, tái xếp hạng bằng Cross-Encoder.
  \item \code{FaithfulnessEvaluator}: Đánh giá độ trung thực của câu trả lời dựa trên việc đối chiếu từng khẳng định nguyên tử với ngữ cảnh trích xuất.
  \item \code{GraphEntityRetriever}: Mở rộng tri thức cục bộ theo đồ thị liên kết $k$-hop.
  \item \code{TableAwareChunker}: Bộ chia đoạn bảng biểu bảo toàn header cho RAG doanh nghiệp:
  \item \code{ColBERTMaxSimEngine}: Động cơ chấm điểm Late Interaction trên ma trận véc-tơ token đa chiều (được kiểm thử trong test case \code{test_colbert_maxsim_numerical_correctness}).
  \item \code{ContextualChunkAugmenter}: Bộ tiền xử lý ngữ cảnh hóa Contextual Retrieval gắn thẻ XML tài liệu (được kiểm thử trong test case \code{test_contextual_chunk_augmenter}).
\end{enumerate}

\begin{lstlisting}[language=Python, caption={Bộ phân đoạn bảng biểu nhận biết cấu trúc TableAwareChunker.}]
class TableAwareChunker:
    """Bo chia doan van ban Markdown co kha nang bao toan Table Headers."""
    def __init__(self, max_table_rows: int = 5):
        self.max_table_rows = max_table_rows

    def chunk_markdown(self, text: str) -> List[str]:
        lines = text.strip().split("\n")
        chunks = []
        in_table = False
        table_headers = []
        current_rows = []

        for line in lines:
            stripped = line.strip()
            if stripped.startswith("|") and stripped.endswith("|"):
                if not in_table:
                    in_table = True
                    table_headers = [line]
                elif len(table_headers) == 1 and ("---" in stripped):
                    table_headers.append(line)
                else:
                    current_rows.append(line)
                    if len(current_rows) >= self.max_table_rows:
                        chunk_content = "\n".join(table_headers + current_rows)
                        chunks.append(chunk_content)
                        current_rows = []
            else:
                if in_table:
                    if current_rows:
                        chunk_content = "\n".join(table_headers + current_rows)
                        chunks.append(chunk_content)
                        current_rows = []
                    in_table = False
                    table_headers = []
                if stripped:
                    chunks.append(line)

        if in_table and current_rows:
            chunks.append("\n".join(table_headers + current_rows))

        return chunks
\end{lstlisting}

\subsection{Lời Giải Nhóm D: Phân Tích Thất Bại Thực Tế (Post-Mortem Analysis)}

\begin{debugbox}[Báo Cáo Sự Cố 1: Xung Đột Phiên Bản Tài Liệu \& Lỗi Trôi Thời Gian]
\textbf{1. Tóm tắt sự cố}: Hệ thống RAG trả về mức lãi suất cũ $7.2\%$ thay vì lãi suất mới $5.5\%$.
\textbf{2. Phân tích nguyên nhân gốc rễ}: Mô hình nhúng véc-tơ không nhận thức về dòng thời gian (Time-agnostic). Tài liệu cũ diễn đạt dài và chứa nhiều từ khóa trùng với câu hỏi hơn tài liệu mới.
\textbf{3. Giải pháp khắc phục}:
\begin{itemize}
  \item Gắn metadata cứng: \code{effective_date}, \code{status: "active"}.
  \item Áp dụng lọc Metadata Pre-filtering trước khi tính khoảng cách véc-tơ.
  \item Thiết lập cơ chế tự động xóa hoặc đánh dấu \code{deprecated} các tài liệu đã hết hiệu lực.
\end{itemize}
\end{debugbox}

\begin{debugbox}[Báo Cáo Sự Cố 2: Lỗi Chunking Cắt Ngang Bảng Biểu Tài Chính]
\textbf{1. Tóm tắt sự cố}: Hệ thống RAG nhầm lẫn số liệu doanh thu năm 2024 thành năm 2025 do bộ chia đoạn cắt ngang giữa tiêu đề cột và hàng số liệu.
\textbf{2. Phân tích nguyên nhân gốc rễ}:
Thuật toán phân đoạn văn bản thông thường (Character-based hoặc Token-based Text Splitter) là thuật toán \textit{mù cấu trúc} (\textit{Structure-agnostic}). Nó chỉ đếm đủ 512 ký tự rồi bẻ đôi đoạn văn bản, bất kể vị trí cắt đang nằm giữa một bảng Markdown hay một câu lệnh code.
\textbf{3. Giải pháp khắc phục (Table-Aware Semantic Chunking)}:
\begin{enumerate}
  \item Sử dụng bộ phân tích cú pháp cấu trúc (như Unstructured hoặc LlamaParse) để nhận diện đối tượng bảng biểu (Table Object).
  \item \textbf{Cơ chế Table Header Injection}: Khi một bảng biểu quá dài buộc phải chia nhỏ, mỗi chunk con bắt buộc phải sao chép lại toàn bộ hàng tiêu đề cột (\code{Table Headers}) và tên bảng ở ngay dòng đầu tiên.
  \item Chuyển đổi bảng biểu thành định dạng mô tả tự nhiên: Sử dụng LLM tóm tắt từng hàng thành một câu hoàn chỉnh (ví dụ: \textit{``Trong năm 2025, doanh thu mảng bán lẻ đạt 1.850 tỷ đồng''}) trước khi nhúng véc-tơ.
\end{enumerate}
\end{debugbox}

\begin{checkpointbox}[Tiêu Chuẩn Đạt Yêu Cầu Về Kỹ Nghệ RAG Doanh Nghiệp]
Một kỹ sư AI hoàn thành xuất sắc Chương 8 khi:
\begin{itemize}
  \item Phân tích cặn kẽ 4 điểm nghẽn vật lý của Naive RAG trong thực tế triển khai.
  \item Nắm vững cơ sở toán học của IVFFlat, HNSW, BM25, và thuật toán dung hợp thứ hạng RRF.
  \item Hiểu rõ giới hạn của Vector RAG trước bài toán Global Sensemaking và kiến trúc GraphRAG phân cấp.
  \item Nắm vững quy trình phản hồi tự thích ứng của Corrective RAG (CRAG) và cách đo lường RAG Triad.
\end{itemize}
\end{checkpointbox}
